% theory/ParabolicDiffusion/GEAR_FVPM_Diffusion_Theory.tex
% Master document. Preamble and front matter only -- the actual physics
% content lives in parabolic.tex, hyperbolic.tex, and appendices.tex,
% pulled in below with \input. Edit those files directly; each covers one
% topic and keeps its own focused diff. Regenerate the PDF with ./run.sh.

\documentclass[11pt, a4paper, parskip]{scrartcl}
\usepackage[margin=2.7cm]{geometry}
\usepackage{amsmath}
\usepackage{amssymb}
\usepackage{amsthm}
\usepackage{graphicx}
\usepackage[utf8]{inputenc} % Required for Unicode characters
\usepackage{xcolor}         % Required for coloring text
\usepackage{enumitem}       % For customizing list environments
\usepackage{physics}        % For common physics notations like dv, grad, div, etc.
\usepackage{booktabs}       % For clean tables
\usepackage{seqsplit}
\usepackage[hidelinks]{hyperref}

\usepackage[]{natbib}
\bibliographystyle{apalike}

% \code{} is \texttt{} that may break mid-identifier/mid-path across lines
% (long function names, file paths) instead of overflowing the margin,
% used for long \texttt content in running text. Left as plain \texttt{}
% inside captions/section titles, where \seqsplit's \futurelet-based
% scanning cannot survive LaTeX's full expansion of moving arguments.
\newcommand{\code}[1]{\texttt{\seqsplit{#1}}}

% KOMA-Script disables the classic \rm/\bf/\it/... font-switch commands
% by default; the numbered sections use them freely for math-mode units
% ($h_{\rm phys}$, $c_{\rm hyp}$, ...), matching the old article.cls
% setup, so restore the standard LaTeX2e old-font-command behaviour here
% rather than editing every occurrence.
\makeatletter
\DeclareOldFontCommand{\rm}{\normalfont\rmfamily}{\mathrm}
\DeclareOldFontCommand{\sf}{\normalfont\sffamily}{\mathsf}
\DeclareOldFontCommand{\tt}{\normalfont\ttfamily}{\mathtt}
\DeclareOldFontCommand{\bf}{\normalfont\bfseries}{\mathbf}
\DeclareOldFontCommand{\it}{\normalfont\itshape}{\mathit}
\DeclareOldFontCommand{\sl}{\normalfont\slshape}{\@nomath\sl}
\DeclareOldFontCommand{\sc}{\normalfont\scshape}{\@nomath\sc}
\makeatother

\setkomafont{title}{\rmfamily\bfseries\boldmath}
\addtokomafont{section}{\rmfamily\bfseries\boldmath}
\addtokomafont{subsection}{\rmfamily\bfseries\boldmath}
\addtokomafont{subsubsection}{\rmfamily\bfseries\boldmath}

\theoremstyle{plain}
\newtheorem{proposition}{Proposition}
\newtheorem{lemma}{Lemma}
\newtheorem{corollary}{Corollary}
\theoremstyle{definition}
\newtheorem{definition}{Definition}
\theoremstyle{remark}
\newtheorem{remark}{Remark}
\newtheorem{example}{Example}

% ---------------------------------------------------------------------------
% Physics macros, shared between the parabolic and hyperbolic parts. Names
% are chosen to mirror the GEAR\_FVPM\_DIFFUSION source (see the
% Nomenclature appendix for the symbol <-> code-identifier map).
% ---------------------------------------------------------------------------
\newcommand{\Zfrac}{Z}                                  % metal mass fraction
\newcommand{\rhoZ}{\rho \Zfrac}                          % metal mass density
\newcommand{\driver}{q}                                  % diffusion driver
\newcommand{\Kmat}{\mathbf{K}}                            % diffusion tensor
\newcommand{\Kstar}{\mathbf{K}^{\star}}                   % interface diffusion tensor
\newcommand{\Shear}{\mathbf{S}}                           % shear tensor
\newcommand{\Splus}{\mathbf{S}_{+}}                       % regularised (PSD) shear tensor
\newcommand{\Fdiff}{\mathbf{F}}                           % diffusive flux (algebraic in Part I, dynamical from Part II on)
\newcommand{\kappaC}{\kappa_{0}}                          % chem\_data->diffusion\_coefficient
\newcommand{\kappaP}{\kappa}                              % per-particle chemistry\_data.kappa
\newcommand{\CCFL}{C_{\mathrm{CFL}}}                      % chem\_data->C\_CFL\_chemistry
\newcommand{\psiHLL}{\psi}                                % hll\_riemann\_solver\_psi
\newcommand{\epsHLL}{\epsilon}                            % hll\_riemann\_solver\_epsilon
\newcommand{\gammak}{\gamma_{k}}                          % kernel\_gamma
\newcommand{\Aij}{A_{ij}}                                 % interface area
\newcommand{\xij}{\mathbf{x}_{ij}}                        % interface position
\newcommand{\nunit}{\mathbf{n}}                           % interface unit normal
\newcommand{\dx}{\Delta x}                                % GIZMO particle size
\newcommand{\dt}{\Delta t}                                % explicit timestep
\newcommand{\Uvec}{\mathbf{U}}                            % hyperbolic state vector
\newcommand{\Gflux}[1]{\mathbf{G}^{(#1)}}                 % flux tensor, direction-#1 component
\newcommand{\chyp}{c_{\rm hyp}}                           % hyperbolic diffusion sound speed
\newcommand{\taurelax}{\tau}                              % relaxation time
\newcommand{\Deff}{\mathbf{D}}                            % effective diffusivity tensor, D = K*(q/U)
\newcommand{\expr}{\mathcal{E}}                           % chemistry_timesteps.h's gradient-suppression "expression"

\title{Theoretical Foundations of the GEAR FVPM Metal-Diffusion Scheme}
\subtitle{Theory document: parabolic (Fickian) and hyperbolic
(Cattaneo-relaxation) formulations}
\author{Darwin Roduit}
\date{\today}

\begin{document}

\maketitle

This is the theoretical and numerical reference for both variants of the
GEAR FVPM metal-diffusion chemistry scheme
(\code{src/chemistry/GEAR\_FVPM\_DIFFUSION/}): the parabolic (Fickian)
formulation, configured with
\code{--with-chemistry=GEAR-FVPM-DIFFUSION\_N} (Part~\ref{part:parabolic}),
and the hyperbolic (Cattaneo-relaxation) formulation, configured with
\code{--with-chemistry=GEAR-FVPM-HYPERBOLIC-DIFFUSION\_N}
(\code{src/chemistry/GEAR\_FVPM\_DIFFUSION/hyperbolic/}, from
Part~\ref{part:hyperbolic}). Every equation is cross-referenced to the C
source. User-facing configuration and parameters are documented in the
RTD pages under \code{doc/RTD/source/SubgridModels/GEAR/diffusion/}.
Regenerate this PDF with \code{./run.sh}.

\tableofcontents

% theory/ParabolicDiffusion/parabolic.tex
% Body of the parabolic (Fickian) FVPM diffusion scheme. Input from
% GEAR_FVPM_Diffusion_Theory.tex (the master document); not standalone.

\part{Parabolic (Fickian) Diffusion}
\label{part:parabolic}

\section{Introduction}

\subsection{Purpose and Scope}
\label{sec:parabolic-scope}
This document is a rigorous theoretical and numerical account of the
\emph{parabolic} metal-diffusion solver implemented in
\code{src/chemistry/GEAR\_FVPM\_DIFFUSION} (shared headers) and
\code{src/chemistry/GEAR\_FVPM\_DIFFUSION/parabolic} (mode-specific headers)
of the GEAR Finite Volume Particle Method (FVPM) chemistry module. Every
equation below is transcribed from, and cross-referenced to, the C source
so that the document can be used as a derivation-level companion to the
code rather than a qualitative overview.

The code also contains a second, \emph{hyperbolic} (Cattaneo/Maxwell-type)
formulation of metal diffusion, selected at compile time by the macro
\code{CHEMISTRY\_GEAR\_FVPM\_HYPERBOLIC\_DIFFUSION} and implemented in the
sibling \code{hyperbolic/} subdirectory. That formulation is out of scope
for this part: the parabolic equation solved here is recovered as its
relaxation-time-to-zero limit, and the two share the FVPM geometry, the
diffusion driver $\driver$, and the diffusion tensor $\Kmat$ developed
below, but use different Riemann solvers, conserved-state vectors, and
timestep criteria. The hyperbolic solver's own derivation, characteristic
speed, and exact solutions are covered in Part~\ref{part:hyperbolic}.

\subsection{Overview of the GEAR FVPM Framework}
GEAR's chemistry module rides on top of a meshless Finite Volume Particle
Method (Lagrangian, GIZMO/MFM/MFV-style) hydrodynamics solver. Each particle
$i$ carries an effective volume $V_i$ and a geometric moment matrix
$\mathbf{B}_i$ (SWIFT field \code{p->geometry.matrix\_E}, already inverted)
that is computed once per particle in the hydrodynamics density loop and
consumed, unmodified, by the chemistry module. The diffusion solver
therefore does not own its own spatial discretisation: it reuses the
host FVPM geometry to build kernel-weighted least-squares gradients and
generalised interface areas between neighbouring particles (see
Section~\ref{sec:fvpm-geometry}).

\subsection{Metal Transport in Astrophysical Simulations}
Unresolved turbulent mixing redistributes heavy elements produced by
stars on scales below the hydrodynamic resolution. Because Lagrangian
particle methods do not exchange mass between neighbours by default
(each particle retains its metal budget unless an explicit mixing
operator is added), an explicit sub-grid diffusion term is required to
reproduce a metal mixing of the interstellar and intergalactic medium.
This is the role of the module documented here.

\section{Physical Foundations}
\label{sec:physical-foundations}

\subsection{Governing Differential Equations}

\subsubsection{Conservation of Metal Mass}
Let $\rho$ be the gas density and $\Zfrac$ the mass fraction of a given
metal species (code index \code{metal} $\in [0,$
\code{GEAR\_CHEMISTRY\_ELEMENT\_COUNT}$)$). The transported quantity is the
metal mass density $\rhoZ$. In the absence of sources (star-formation
depletion) or sinks (supernova/stellar-wind enrichment, handled by other
modules), conservation of metal mass advected with the flow and diffused
relative to it reads:
\begin{equation}
  \label{eq:conservation_rho_Z}
  \pdv{(\rhoZ)}{t} + \div (\rhoZ\, \mathbf{v}) = -\div \Fdiff ,
\end{equation}
where $\mathbf{v}$ is the gas velocity and $\Fdiff$ is the diffusive
metal-mass flux. The advective term $\div(\rhoZ\,\mathbf v)$ is handled
implicitly by the fact that $\rhoZ$ is carried as a passive scalar by the
Lagrangian particles (their mass and position already advect with the
flow); the chemistry module solves only the remaining, purely diffusive
part in the frame co-moving with each particle.

\subsubsection{Parabolic Diffusion Equation}
Isolating the diffusive part of Eq.~\eqref{eq:conservation_rho_Z} gives
the parabolic diffusion equation actually integrated by the code
(\code{chemistry\_end\_force}, \code{chemistry.h}):
\begin{equation}
  \label{eq:parabolic_diffusion_rho_Z}
  \pdv{(\rhoZ)}{t} = -\div \Fdiff .
\end{equation}
This equation is discretised in finite-volume (conservative) form: each
particle's metal mass $m_Z = \rhoZ\, V$ is updated by summing the flux
exchanged with every neighbour across the shared interface (see
Section~\ref{sec:flux-computation}).

\subsubsection{Fick's Law and the Diffusion Driver}
\label{sec:ficks-law}
The diffusive flux obeys a generalised (tensorial) Fick's law,
\begin{equation}
  \label{eq:ficks_law}
  \Fdiff = -\Kmat \cdot \grad \driver ,
\end{equation}
where $\Kmat$ is a symmetric, positive semi-definite $3\times3$ diffusion
tensor (Section~\ref{sec:diffusion-tensor}) and $\driver$ is a
\emph{diffusion driver} whose identity is not fixed: depending on the
diffusion mode (Section~\ref{sec:diffusion-modes}) the code sets
$\driver = \rhoZ$ or $\driver = \Zfrac$. This choice is physically
consequential: it changes which PDE is actually solved (see the boxed remark at the end of
Section~\ref{sec:diffusion-modes}).

\subsection{Assumptions and Simplifications}

\textbf{Continuum approximation.}
Equations~\eqref{eq:conservation_rho_Z}--\eqref{eq:ficks_law} presuppose
that the gas and metal fields are continuous on scales larger than the
particle spacing, i.e.\ that the FVPM discretisation resolves the scales
of interest.

\textbf{Operator splitting from other physics.}
Diffusion is solved with operator splitting: hydrodynamics, gravity,
cooling, star formation and feedback are integrated by other SWIFT
modules on their own schedules, and the chemistry module only ever sees
the metal mass field at the start and end of its own sub-step
(Section~\ref{sec:time-integration}).

\textbf{Isotropy versus anisotropy.}
$\Kmat$ can be isotropic ($\Kmat = \kappaP \mathbf I$) or anisotropic
(full symmetric tensor built from the local velocity shear). Both cases
are implemented; see Section~\ref{sec:diffusion-modes}.

\subsection{Metal Transport Parameters}

\subsubsection{Metal Mass Fraction}
$\Zfrac = m_Z / m$, with $m$ the particle's total (gas) mass. Computed by
\code{chemistry\_get\_metal\_mass\_fraction} (\code{chemistry\_getters.h}).
It is dimensionless and, physically, bounded in $[0,1]$; the code enforces
this bound only at specific points (Section~\ref{sec:unphysical}), since
the diffusion solver internally tolerates small transient excursions
below zero to control noise amplification (Section~\ref{sec:gradients}).

\subsubsection{Diffusion Coefficient and Tensor}
\label{sec:diffusion-tensor}
The single global input parameter \code{chem\_data->diffusion\_coefficient}
(denoted $\kappaC$ below) is re-used with three distinct physical
meanings depending on the diffusion mode. Note the resulting unit
change: in \code{isotropic\_constant} mode $\kappaC$ is a true diffusivity
(units of length$^2$/time); in the other two modes it plays the role of a
dimensionless Smagorinsky-like constant $C_S$ multiplying a
resolution-dependent turbulent diffusivity.

\subsubsection{Shear Tensor}
\label{sec:shear-tensor}
For the turbulence-coupled modes, $\Kmat$ depends on the local,
trace-free (compression/expansion removed) velocity shear tensor
$\Shear$, built from a kernel-filtered ("Favre-like") velocity field
$\tilde{\mathbf v}$ rather than the raw particle velocity, in the spirit
of a large-eddy-simulation (LES) closure. Its full construction is given
in Section~\ref{sec:diffusion-modes}.

\subsection{Diffusion Modes}
\label{sec:diffusion-modes}

The behaviour of the solver is controlled by
\code{chem\_data->diffusion\_mode},
which selects one of three closures. In all three cases the diffusion
coefficient is computed once per active particle per force sub-step by
\code{chemistry\_get\_diffusion\_coefficient}
(\code{chemistry\_getters.h}), using the physical (proper, not
comoving) filtered density $\bar\rho = a^{-3}\,$\code{p->chemistry\_data.filtered.rho}
and physical smoothing length $h_{\rm phys} = a\,h$.

\subsubsection{Isotropic Constant Diffusion}
\begin{equation}
  \Kmat = \kappaC \, \mathbf I , \qquad \driver = \rhoZ .
\end{equation}
This is the simplest closure: $\kappaC$ is a user-specified constant
diffusivity, independent of the local flow. The resulting PDE
(Section~\ref{sec:pde-per-mode}) is a linear heat equation for
the metal density itself,
\begin{equation}
  \Fdiff = -\kappaC \,\grad(\rhoZ) ,
\end{equation}
useful primarily as a verification test against analytic diffusion
solutions (e.g.\ the spreading of a Gaussian or delta-function peak; see
the \code{MetalDiffusionOnePeak}/\code{MetalDiffusionTwoPeaks} test
problems shipped alongside this module).

\subsubsection{Isotropic Smagorinsky Model}
\begin{equation}
  \kappaP = \kappaC \, \gammak^{2} \, h_{\rm phys}^{2} \, \bar\rho \,
            \norm{\Shear}_F , \qquad
  \Kmat = \kappaP \, \mathbf I , \qquad \driver = \Zfrac ,
\end{equation}
where $\norm{\Shear}_F = \sqrt{\Shear:\Shear}$ is the Frobenius norm of
the shear tensor of Section~\ref{sec:shear-def} and $\gammak$ is the
kernel's compact-support radius factor (SWIFT's \code{kernel\_gamma},
$O(1$--$2)$ and kernel-choice dependent, so that
$h_{\rm bar} = \gammak\, h$ is the true compact-support radius of the
kernel of smoothing length $h$). This is a standard Smagorinsky
sub-grid-scale closure: the effective diffusivity scales as
(filter length)$^2\times$(local strain rate), exactly as an eddy
viscosity would in an LES momentum closure. The flux is
\begin{equation}
  \Fdiff = -\kappaP \,\grad \Zfrac .
\end{equation}

\subsubsection{Anisotropic Gradient Model}
\begin{equation}
  \kappaP = \kappaC \, \gammak^{2} \, h_{\rm phys}^{2} \, \bar\rho , \qquad
  \Kmat = \kappaP \, \Splus , \qquad \driver = \Zfrac ,
\end{equation}
i.e.\ the same resolution-dependent prefactor as the Smagorinsky mode,
\emph{without} the $\norm{\Shear}_F$ factor, because directional
information is now retained in full through the tensor $\Splus$ rather
than collapsed into a scalar norm (Rennehan-style~\cite{Rennehan2021} gradient/shear-aligned
diffusion). The flux is the full matrix contraction
\begin{equation}
  \Fdiff = -\Kmat \cdot \grad \Zfrac
         = -\kappaP\, \Splus \cdot \grad \Zfrac .
\end{equation}

\subsubsection{Construction of the Shear Tensor and its Regularisation}
\label{sec:shear-def}
The raw (traceless, Hubble-flow-corrected) shear tensor is built in
\code{chemistry\_get\_physical\_shear\_tensor}
(\code{chemistry\_getters.h}) from the symmetric part of the
gradient of the filtered velocity $\tilde{\mathbf v}$ (comoving
gradients, converted to physical units by $a^{-2}$: one factor of
$a^{-1}$ from $\mathbf v_{\rm peculiar} = \mathbf v_{\rm comoving}/a$ and
one from $\grad_{\rm phys} = a^{-1}\grad_{\rm comoving}$):
\begin{equation}
  \label{eq:shear_raw}
  S_{ij} = \frac{1}{2}
           \left(\pdv{\tilde v_i}{x_j} + \pdv{\tilde v_j}{x_i}\right)_{\rm phys}
           + H\,\delta_{ij}
           - \frac{1}{3}\,\mathrm{tr}(\Shear')\,\delta_{ij} ,
\end{equation}
where $H$ is the Hubble parameter (added once to the diagonal, since the
Hubble flow is isotropic and $\Shear$ is symmetric) and the final term
removes the trace of $\Shear' \equiv \tfrac12(\grad\tilde{\mathbf v} +
\grad\tilde{\mathbf v}^T) + H\mathbf I$ \emph{after} the Hubble term has
been added, so that pure compression/expansion (which is not turbulence)
never contributes to the diffusion coefficient. The filtered velocity
$\tilde{\mathbf v} = \overline{\rho\mathbf v}/\bar\rho$ is a
mass(Favre)-weighted, kernel-smoothed velocity evaluated over the full
compact support $h_{\rm bar} = \gammak h$ (wider than the density-defining
smoothing length $h$), accumulated pairwise in
\code{runner\_iact\_chemistry}/\code{runner\_iact\_nonsym\_chemistry}
(\code{chemistry\_iact.h}) with a harmonic-mean density weighting
and a fixed damping prefactor
\code{GEAR\_FVPM\_DIFFUSION\_FILTERING\_SMOOTHING\_FACTOR}$= 0.8$
(\code{chemistry.h}) applied to the accumulated kernel sum before the
particle's own self-contribution $\bar\rho_i \mathbf v_i$ is added back
in \code{chemistry\_end\_density} (\code{chemistry.h}). Filtering
$\mathbf v$ over the wider kernel support before differentiating
suppresses single-particle noise in the shear estimate, consistent with
the LES philosophy underlying the Smagorinsky and anisotropic-gradient
closures.

$\Shear$ as constructed in Eq.~\eqref{eq:shear_raw} is symmetric but not
sign-definite. For the anisotropic mode this is a problem: $\Kmat$ must
be positive semi-definite (PSD), otherwise the flux $\Fdiff = -\Kmat
\cdot\grad\Zfrac$ can point \emph{up} the metallicity gradient in some
direction and the diffusion operator would sharpen rather than smooth
metal contrasts, an unconditionally unstable configuration for an
explicit parabolic solver. \code{chemistry\_regularize\_shear\_tensor}
(\code{chemistry\_getters.h}) fixes this by discarding the
negative-eigenvalue subspace of $\Shear$, following~\cite{Balarac2013}:
\begin{equation}
  \label{eq:S_plus_def}
  \Splus \;=\; \sum_{k=1}^{3} \max(0,\lambda_k)\, \mathbf e_k \otimes \mathbf e_k ,
\end{equation}
where $(\lambda_k, \mathbf e_k)$ are the eigenpairs of the symmetric
matrix $\Shear$ (obtained from a closed-form $3\times3$ symmetric
Jacobi-type diagonalisation, \code{sym\_matrix\_diagonalise\_3x3\_d}).

\begin{proposition}[Positive semi-definiteness of $\Splus$]
\label{prop:psd}
The tensor $\Splus$ defined by Eq.~\eqref{eq:S_plus_def} is symmetric
and positive semi-definite for any symmetric input $\Shear$.
\end{proposition}
\begin{proof}
Symmetry is immediate since $\mathbf e_k\otimes\mathbf e_k$ is symmetric
for every $k$ and the $\lambda_k$ are real (eigenvalues of a symmetric
real matrix). For positive semi-definiteness, take any
$\mathbf u \in \mathbb R^3$ and expand it in the orthonormal eigenbasis
$\{\mathbf e_k\}$ of $\Shear$ (orthonormality follows from the spectral
theorem applied to the symmetric matrix $\Shear$):
\begin{equation}
  \mathbf u^T \Splus\, \mathbf u
  = \sum_{k=1}^3 \max(0,\lambda_k)\, \mathbf u^T (\mathbf e_k \otimes \mathbf e_k) \mathbf u
  = \sum_{k=1}^3 \max(0,\lambda_k)\, (\mathbf e_k \cdot \mathbf u)^2 .
\end{equation}
Each term in the sum is a non-negative scalar $\max(0,\lambda_k)$ times a
square, hence non-negative; the sum of non-negative terms is
non-negative, so $\mathbf u^T \Splus\,\mathbf u \geq 0$ for all
$\mathbf u$, i.e.\ $\Splus \succeq 0$.
\end{proof}

Because $\Kmat = \kappaP \Splus$ with $\kappaP \geq 0$, Proposition
\ref{prop:psd} guarantees $\Kmat \succeq 0$ in the anisotropic mode by
construction, for every particle, at every timestep, independent of the
local flow topology.

\subsubsection{Different Modes Solve Different PDEs}
\label{sec:pde-per-mode}
\begin{quote}
Substituting the driver choice of each mode into
Eq.~\eqref{eq:parabolic_diffusion_rho_Z} via Eq.~\eqref{eq:ficks_law}
gives three different governing equations, not merely three
parametrisations of the same one:
\begin{align}
  \text{isotropic\_constant:} &\quad
    \pdv{(\rhoZ)}{t} = \div\!\big(\kappaC\, \grad(\rhoZ)\big) ,
    \label{eq:pde_const}\\
  \text{isotropic\_smagorinsky:} &\quad
    \pdv{(\rhoZ)}{t} = \div\!\big(\kappaP(\mathbf x,t)\, \grad \Zfrac\big) ,
    \label{eq:pde_smag}\\
  \text{anisotropic\_gradient:} &\quad
    \pdv{(\rhoZ)}{t} = \div\!\big(\Kmat(\mathbf x,t) \cdot \grad \Zfrac\big) .
    \label{eq:pde_aniso}
\end{align}
Equations~\eqref{eq:pde_smag}--\eqref{eq:pde_aniso} diffuse the
\emph{intensive} quantity $\Zfrac$, with the local density folded into
the (state-dependent) coefficient $\kappaP \propto \bar\rho$; this is the
standard sub-grid metal-mixing closure used in the literature
(e.g.\ Shen, Wadsley \& Stinson~\citeyear{ShenWadsleyStinson2010}; Escala et al.~\citeyear{Escala2018}), chosen
precisely because diffusing $\rhoZ$ directly (as in
Eq.~\eqref{eq:pde_const}) would let density gradients alone drive
spurious metal currents even in a chemically homogeneous ($\grad Z = 0$)
medium. Equation~\eqref{eq:pde_const} does not have this property and is
intended for controlled test problems rather than production
astrophysical runs.
\end{quote}

\section{FVPM Spatial Discretisation}
\label{sec:fvpm-geometry}

\subsection{Meshless Particle Representation}
The computational domain is discretised by particles, each carrying an
effective volume $V_i$ (\code{p->geometry.volume}) computed by the
underlying FVPM solver from a kernel partition of unity.
Interactions occur between particles found by the same neighbour search
(compact-support kernel with smoothing length $h$) used by the
hydrodynamics solver; the chemistry module performs no independent
neighbour finding.

\subsection{Least-Squares Gradient Estimator}
\label{sec:gradient-estimator}
For every scalar field $Q$ carried by the particles (metal mass
fraction, metal density, velocity components), the module computes a
kernel-weighted, matrix-corrected gradient estimate
(\code{chemistry\_gradients\_collect}, \code{chemistry\_gradients.h}):
\begin{equation}
  \label{eq:gradient_estimator}
  (\grad Q)_i \;=\; \sum_{j} (Q_i - Q_j)\, \tilde{\boldsymbol\psi}_i(\mathbf x_j) ,
  \qquad
  \tilde{\boldsymbol\psi}_i(\mathbf x_j) = W_{ij}\, \mathbf B_i \cdot \mathbf{dx}_{ij} ,
\end{equation}
with $\mathbf{dx}_{ij} \equiv \mathbf x_i - \mathbf x_j$ (the code's
sign convention), $W_{ij} \equiv W(|\mathbf{dx}_{ij}|, h_i)$ the kernel
value, and $\mathbf B_i = \mathbf E_i^{-1}$ the inverse of the particle's
geometric moment matrix (\code{p->geometry.matrix\_E}), computed and
maintained by the host hydrodynamics module using the GIZMO second-moment
normalisation of Lanson \& Vila~\citeyear{LansonVila2008} and
Hopkins~\citeyear{Hopkins2015},
\begin{equation}
  \label{eq:E_matrix_def}
  E_i^{kl} = \sum_j W_{ij}\; \mathrm{dx}_{ij}^{k}\, \mathrm{dx}_{ij}^{l} .
\end{equation}
This is Hopkins~\citeyear{Hopkins2015}'s moment matrix (their eqs.~6--7
and 14) with the particle's own partition-of-unity normalisation
already cancelled out; the code accumulates it directly in
\code{matrix\_E} (\code{fvpm\_accumulate\_geometry\_and\_matrix},
\code{fvpm\_geometry.h}) without computing that normalisation
separately.

\begin{proposition}[Linear consistency of the gradient estimator]
\label{prop:gradient-consistency}
If a scalar field is exactly linear in a neighbourhood of particle $i$,
$Q(\mathbf x) = Q_i + \mathbf G\cdot(\mathbf x - \mathbf x_i)$ for a
constant vector $\mathbf G$, then the estimator of
Eq.~\eqref{eq:gradient_estimator} reproduces the true gradient exactly:
$(\grad Q)_i = \mathbf G$. Proof in
Appendix~\ref{app:gradient-consistency-proof}.
\end{proposition}

Whenever $\mathbf E_i$ is ill-conditioned (near-singular particle
distribution, flagged by \code{fvpm\_part\_geometry\_well\_behaved(p)}
returning false when the stored condition number exceeds
\code{const\_gizmo\_max\_condition\_number}~$=100$), the matrix inversion is
abandoned and the code falls back to the plain (first-order, kernel
gradient) SPH estimator
\begin{equation}
  \tilde{\boldsymbol\psi}_i(\mathbf x_j) = -\frac{W'_{ij}}{r_{ij}}\, \mathbf{dx}_{ij} ,
\end{equation}
with $W'_{ij}$ the kernel derivative and $r_{ij} = |\mathbf{dx}_{ij}|$
(\code{chemistry\_gradients.h}). After accumulation over all
neighbours, gradients are normalised by $h^{-d}$ ($d$ the spatial
dimension) in \code{chemistry\_gradients\_finalise}
(\code{chemistry\_gradients.h}), matching the kernel's own
normalisation convention.

\subsection{Interface Geometry: Generalised Face Area}
\label{sec:face-area}
Fluxes are exchanged across the interface between every
interacting pair $(i,j)$. Following Hopkins (2017, eq.~7), the
generalised interface area vector is
\begin{equation}
  \label{eq:face_area}
  \mathbf A_{ij} = -X_i\, W_{ij}\, h_i^{-d}\, \big(\mathbf B_i \cdot \mathbf{dx}_{ij}\big)
                   -X_j\, W_{ji}\, h_j^{-d}\, \big(\mathbf B_j \cdot \mathbf{dx}_{ij}\big) ,
\end{equation}
with $X_i = V_i$, $X_j = V_j$ in the well-behaved case. If the two
particle volumes differ strongly,
$|V_i-V_j|/\min(V_i,V_j) > 1.5\, d$ (the \code{GIZMO\_VOLUME\_CORRECTION}
branch, \code{chemistry\_iact.h}), both weights are replaced by
the smoothing-length-weighted blend
$X_i = X_j = (V_i h_j + V_j h_i)/(h_i+h_j)$
to avoid numerically unstable cancellations in
Eq.~\eqref{eq:face_area} across sharp density (and hence volume)
contrasts. If either particle's geometry matrix is ill-conditioned, the
code instead falls back to the isotropic SPH face-area estimate
\begin{equation}
  A_{ij}^{\rm SPH} = -\Big(h_i^{-(d+1)} V_i^2 W'_{ij} + h_j^{-(d+1)} V_j^2 W'_{ji}\Big)/r_{ij} ,
  \qquad \mathbf A_{ij} = A_{ij}^{\rm SPH}\, \mathbf{dx}_{ij} .
\end{equation}
The interface normal is $\nunit = \mathbf A_{ij}/\Aij$ with
$\Aij = |\mathbf A_{ij}|$
(\code{chemistry\_iact.h}).

\subsection{Interface Position and Interface-Frame Boost}
\label{sec:interface-position}
The effective interface point, relative to particle $i$, is (Hopkins
2017, eq.~8)
\begin{equation}
  \label{eq:interface_position}
  \xij = x_{\rm fac}\, \mathbf{dx}_{ij} , \qquad
  x_{\rm fac} = -\frac{h_i}{h_i+h_j} ,
\end{equation}
so that the resolved interface sits at
$\mathbf x_i + \xij = (h_j\, \mathbf x_i + h_i\, \mathbf x_j)/(h_i+h_j)$:
a smoothing-length-weighted point between the two particles, pulled
towards the particle with the \emph{smaller} smoothing length. The corresponding interface (fluid) velocity used to
boost the hydrodynamic primitive variables into the interface frame
before the Riemann problem is solved is (Hopkins~\citeyear{Hopkins2017diffusion}, eq.~9)
\begin{equation}
  \mathbf v_{ij} = \mathbf v_i + (\mathbf v_i - \mathbf v_j)\, x_{\rm fac}
  \qquad \text{(\code{chemistry\_iact.h})} .
\end{equation}
This position and boost are computed once per pair and reused
identically for the hydrodynamic reconstruction and for the diffusion
Riemann problem, so that both solvers agree on where and in which frame
the interface lives.

\section{Gradient Reconstruction and Slope Limiting}
\label{sec:gradients}

\subsection{Cell-Wise Extremum-Preserving Limiter}
Before any face reconstruction, gradients are scaled down so
that a linear extrapolation to the farthest interacting neighbour cannot
create a value outside the range spanned by the particle and all its
neighbours (Hopkins~\citeyear{Hopkins2015} formalism,
\code{chemistry\_slope\_limiters\_cell.h}). For a scalar $Q$ with cell value
$Q_i$, neighbour extrema $Q_{\min}, Q_{\max}$ and maximal neighbour
distance $r_{\max}$ (\code{chi->limiter.maxr}), the limiter factor
$\alpha\in[0,1]$ is (\code{chemistry\_slope\_limit\_quantity},
\code{chemistry\_slope\_limiters\_cell.h})
\begin{align}
  g &= r_{\max}\, |\grad Q|_i , \\
  \beta &= \max\!\Big(\beta_{\min},\ \beta_{\max}\min\!\Big(1,\ \tfrac{N_{\rm cond}}{\kappa(\mathbf E_i)}\Big)\Big) ,
  \label{eq:cell_limiter_beta} \\
  \alpha_0 &= \min\!\left(1,\ \beta \, \min\!\left(\frac{Q_{\max}-Q_i}{g},\ \frac{Q_i - Q_{\min}}{g}\right)\right) ,
\end{align}
with $\beta_{\min}=1$, $\beta_{\max}=2$
(\code{GIZMO\_SLOPE\_LIMITER\_BETA\_MIN}/\code{MAX}) and $N_{\rm cond}=100$
(\code{const\_gizmo\_max\_condition\_number}) a reference condition number:
as the particle's geometry matrix approaches singularity
($\kappa(\mathbf E_i)\to N_{\rm cond}$), $\beta\to\beta_{\min}$ and the
limiter becomes maximally conservative; in a well-conditioned
configuration $\beta\to\beta_{\max}=2$, the standard Hopkins (2015)
value.

For quantities that must remain non-negative under extrapolation
(the diffusion drivers $\rhoZ$ and $\Zfrac$, invoked with
\code{pos\_preserve=1}), a second, tighter bound is applied
(\code{chemistry\_slope\_limiters\_cell.h}):
\begin{align}
  f_{\rm shoot} &= \min\!\big(Q_i - Q_{\min} + \tau_{\rm shoot}(Q_{\max}-Q_i),\ Q_{\max}-Q_i\big) , \\
  Q_{\min}^{\rm floor} &= \max\!\left(f_{\rm shoot},\ \tfrac12(Q_i+Q_{\min})\right) , \\
  \alpha_{\rm pos} &= \frac{Q_i - Q_{\min}^{\rm floor}}{g} , \\
  \alpha &= \min(\alpha_0,\ \alpha_{\rm pos}) ,
\end{align}
where $\tau_{\rm shoot} = 0.2$
(\code{GEAR\_FVPM\_DIFF\_CELL\_LIMITER\_SHOOT\_TOLERANGE}) is an allowed
fractional overshoot that keeps the limiter from stalling the diffusion
front entirely in regions with strong local extrema. If
$Q_{\min}=Q_{\max}=Q_i$ (an isolated local extremum, e.g.\ a lone
enriched particle surrounded by pristine gas) the gradient is forced to
zero exactly ($\alpha=0$), since Eq.~\eqref{eq:cell_limiter_beta}, above,
would otherwise return a numerically enormous $\alpha$ from a
$0/0$-adjacent limit.

\textbf{Important asymmetry.} The generic cell limiter
\code{chemistry\_slope\_limit\_cell} (called once at the end of the
gradient loop, \code{chemistry\_gradients.h}) is applied to the
velocity gradients but \emph{deliberately not} to the $\Zfrac$/$\rhoZ$
gradients: cell-limiting the
diffusion driver's gradient would permanently reduce the magnitude of
$\grad\driver$ used everywhere it appears in
Eqs.~\eqref{eq:pde_const}--\eqref{eq:pde_aniso}, i.e.\ it would silently
alter the PDE being solved rather than merely stabilise its
discretisation. Instead, the driver gradients are cell-limited
\emph{locally, on the fly}, immediately before each face reconstruction,
inside \code{chemistry\_gradients\_predict}
(parabolic-mode state reconstruction, see below) and
\code{chemistry\_gradients\_predict\_Z}
(\code{chemistry\_gradients.h}), using the same
$\alpha$-computation above but re-evaluated at every pairwise
interaction. This keeps the corrected gradient reserved only for the
extrapolation actually being performed, without corrupting the
particle's stored gradient field for later use in the flux formula
(Section~\ref{sec:flux-computation}), which needs the \emph{uncorrected}
gradient.

\subsection{Face-Wise (GIZMO) Limiter}
\label{sec:face-limiter}
After the cell-limited gradient is linearly extrapolated to the
interface point $\xij$ (Eq.~\eqref{eq:interface_position}), giving a
provisional reconstructed value $\phi_{\rm mid}^{0} = \phi_i + \grad\phi_i
\cdot \xij$, a second limiting pass constrains the extrapolated pair
$(\phi_{\rm mid,i},\phi_{\rm mid,j})$ using both particle values, not
just particle $i$'s neighbourhood extrema
(\code{chemistry\_slope\_limit\_face\_quantity\_double},
\code{chemistry\_slope\_limiters\_face.h}):
\begin{align}
  \delta_1 &= \tfrac12 |\phi_i - \phi_j| , \qquad \delta_2 = \tfrac14 |\phi_i-\phi_j| , \\
  \bar\phi &= \phi_i + \frac{|\xij|}{r_{ij}}(\phi_j - \phi_i) , \\
  \phi_{+} &=
    \begin{cases}
      \phi_{\max} + \delta_1 & \text{if } \mathrm{sign}(\phi_{\max}+\delta_1) = \mathrm{sign}(\phi_{\max}) \\
      \phi_{\max}/\big(1+\delta_1/|\phi_{\max}|\big) & \text{otherwise}
    \end{cases} \\
  \phi_{-} &=
    \begin{cases}
      \phi_{\min} - \delta_1 & \text{if } \mathrm{sign}(\phi_{\min}-\delta_1) = \mathrm{sign}(\phi_{\min}) \\
      \phi_{\min}/\big(1+\delta_1/|\phi_{\min}|\big) & \text{otherwise}
    \end{cases}
\end{align}
with $\phi_{\min}=\min(\phi_i,\phi_j)$, $\phi_{\max}=\max(\phi_i,\phi_j)$.
The final, limited extrapolated value is
\begin{equation}
  \phi_{\rm mid} =
  \begin{cases}
    \max\!\big(\phi_-,\ \min(\bar\phi+\delta_2,\ \phi_{\rm mid}^{0})\big) & \text{if } \phi_i < \phi_j \\
    \min\!\big(\phi_+,\ \max(\bar\phi-\delta_2,\ \phi_{\rm mid}^{0})\big) & \text{otherwise,}
  \end{cases}
\end{equation}
and the returned, limited increment is $\phi_{\rm mid} - \phi_i$. The
sign-aware construction of $\phi_\pm$ avoids letting the bound cross zero
when $\phi_{\min}$ or $\phi_{\max}$ is itself close to zero, which
matters for a strictly non-negative field like $\Zfrac$ or $\rhoZ$.
Optionally (used for the hyperbolic scheme only, not here) the result can
be clamped to non-negative values by $\phi_{\rm mid}\to\max(0,\phi_{\rm
mid})$. For parabolic diffusion this exact limiter is applied
identically to $\Fdiff$'s driver reconstruction via
\code{chemistry\_slope\_limit\_face}
(\code{parabolic/chemistry\_slope\_limiter\_face.h}), which simply
calls the routine above once for each side of the interface ($i\to j$
and $j\to i$).

\subsection{State Reconstruction at the Interface}
\label{sec:state-reconstruction}
The full per-pair, per-species reconstruction pipeline,
\code{chemistry\_gradients\_predict}
(\code{parabolic/chemistry\_gradients.h}), proceeds as follows for
the metal-density state $U \equiv \rhoZ$:
\begin{enumerate}
  \item Start from the comoving metal densities $U_i, U_j$
        (\code{chemistry\_get\_comoving\_metal\_density}).
  \item Cell-limit the stored $\grad(\rhoZ)$ gradients of $i$ and $j$
        \emph{on the fly}, with the positivity-preserving branch enabled
        (Section~\ref{sec:gradients}).
  \item Linearly extrapolate to the interface point $\xij$
        (Section~\ref{sec:interface-position}):
        $\delta U_i = \grad(\rhoZ)_i \cdot \xij$,
        $\delta U_j = \grad(\rhoZ)_j \cdot \xij^{(j)}$, with
        $\xij^{(j)} = \xij + \mathbf{dx}_{ij}$ the same physical point
        expressed relative to particle $j$.
  \item Apply the face-wise limiter of Section~\ref{sec:face-limiter} to
        $(\delta U_i,\delta U_j)$ jointly, then set
        $U_i \mathrel{+}= \delta U_i$, $U_j \mathrel{+}= \delta U_j$.
  \item Correct unphysical excursions
        (\code{chemistry\_gradients\_correct\_unphysical\_states},
        \code{parabolic/chemistry\_gradients.h}; see
        Section~\ref{sec:unphysical} for the full logic): if the
        extrapolation would push the implied metal mass negative or
        above the gas mass, the offending state is reverted to its
        pre-extrapolation (0th-order) value.
  \item \textbf{Final safety clamp}: irrespective of the above,
        $U_i \gets \max(0, U_i)$ and $U_j \gets \max(0, U_j)$
        unconditionally, since the Riemann solver
        (Section~\ref{sec:flux-computation}) is not robust to residual
        negative densities surviving floating-point round-off in the
        limiter chain.
  \item Convert $U_i, U_j$ from comoving to physical density,
        $U \to a^{-3} U$.
\end{enumerate}
When the diffusion driver is $\Zfrac$ rather than $\rhoZ$ (Smagorinsky
and anisotropic modes), the analogous routine
\code{chemistry\_gradients\_predict\_Z}
(\code{chemistry\_gradients.h}) performs the same steps on the
metal mass fraction gradients, followed by a hard clamp of the
reconstructed metallicity to $[0,1]$
(\code{chemistry\_check\_unphysical\_metallicity},
Section~\ref{sec:unphysical}).

\section{Flux Computation}
\label{sec:flux-computation}

\subsection{Wave-Speed Estimation}
The Riemann solver used here (\code{chemistry\_riemann\_solver\_hopkins2017\_HLL},
\code{parabolic/chemistry\_riemann\_HLL.h}) is explicitly noted in
the source as valid \emph{for parabolic diffusion only}: it degrades
for the hyperbolic system when the relaxation time is large. Its wave
speeds are not diffusion signal speeds but ordinary \emph{hydrodynamic}
speeds evaluated at the interface: with $u_L, u_R$ the particle
velocities projected onto the interface normal $\nunit$ and
$c_{s,L}, c_{s,R}$ the (physical) hydrodynamic sound speeds,
\begin{equation}
  \lambda_+ = |u_R - u_L| + \max(c_{s,L}, c_{s,R}) , \qquad \lambda_- = -\lambda_+ .
\end{equation}
This choice guarantees the diffusion Riemann fan is at
least as wide as the hydrodynamic one, so that the diffusive update never
outruns the CFL condition already enforced by the hydro solver on the
same pair. If $\lambda_+-\lambda_- $ falls below a tolerance
$10^{-8}$ (\code{GEAR\_FVPM\_DIFF\_WAVESPEED\_ESTIMATE\_DIFFERENCE\_TOLERANCE}),
both particles are (numerically) at rest relative to the interface with
no local sound speed and the flux is set to zero to avoid a division by
a near-zero denominator.

\subsection{Interface Diffusion Tensor}
The tensor entering the flux is a simple arithmetic mean of the two
particles' physical diffusion tensors,
\begin{equation}
  \Kstar = \tfrac12\big(\Kmat_i + \Kmat_j\big)
  \qquad(\code{chemistry\_riemann\_compute\_K\_star}) .
\end{equation}
If $\norm{\Kstar}_F = 0$ (e.g.\ at the very first timestep, before
$\bar\rho$ has been computed and $\kappaP$ is still zero) the flux is
set to zero identically.

\subsection{The Hopkins~\citeyear{Hopkins2017diffusion} HLL-Type Flux}
\label{sec:hll-flux}
Let $U \equiv \rhoZ$ (the conserved state) and recall $\driver$ is the
diffusion driver, reconstructed to the interface as $\driver_L,
\driver_R$ (Section~\ref{sec:state-reconstruction}), with
$U^\star = \tfrac12(U_L+U_R)$, $\driver^\star = \tfrac12(\driver_L +
\driver_R)$ and $\grad\driver^\star = \tfrac12(\grad\driver_L +
\grad\driver_R)$ (particle-centred, \emph{not} interface-extrapolated,
gradients of the two neighbours).

The flux of each particle's own local, first-order-reconstructed
diffusive flux (Eq.~\eqref{eq:ficks_law}, evaluated with each particle's
own stored gradient) is projected onto the interface normal,
\begin{equation}
  \mathrm{Flux}_{L} = \Fdiff_L\cdot\nunit , \qquad \mathrm{Flux}_R = \Fdiff_R\cdot\nunit .
\end{equation}
The two standard HLL building blocks are then formed,
\begin{equation}
  F_2 = \frac{\lambda_+\,\mathrm{Flux}_L - \lambda_-\,\mathrm{Flux}_R}{\lambda_+-\lambda_-} ,
  \qquad
  F_U = \frac{\lambda_+\lambda_- \,(U_R-U_L)}{\lambda_+-\lambda_-} ,
\end{equation}
i.e.\ $F_2$ is the upwind-blended transport term and $F_U$ is the
standard HLL numerical-dissipation term. Unlike a textbook HLL solver,
which would simply return $F_2 + F_U$, this solver blends the two with a
\emph{minmod} limiter and an \emph{artificial diffusivity} $\alpha$
(Hopkins~\citeyear{Hopkins2017diffusion}, eqs.~7--11) that caps the numerical dissipation at (a
fraction of) the physically prescribed diffusivity:
\begin{equation}
  \label{eq:flux_hll}
  F_{\rm HLL} = \mathrm{minmod}\Big((1+\psiHLL)\, F_2,\ F_2 + \alpha\, F_U\Big) ,
\end{equation}
with the two-argument minmod of Eq.~\eqref{eq:minmod_def} below and
$\psiHLL = $ \code{chem\_data->hll\_riemann\_solver\_psi} a user-tunable
parameter bounding how much extra transport-term weight is admitted
before the antidiffusive correction is clipped.

\subsubsection{The Artificial Diffusivity $\alpha$}
Computed in \code{chemistry\_riemann\_compute\_alpha}
(\code{chemistry\_riemann\_utils.h}):
\begin{align}
  c_{\rm fast}^\star &= \tfrac12(c_{s,L}+c_{s,R}) , \qquad
  v_{\rm HLL} = \tfrac12|u_R-u_L| + c_{\rm fast}^\star , \\
  r &= \frac{v_{\rm HLL}\, |\mathbf{dx}_{ij}|\, |\driver^\star|}{\norm{\Kstar}_F \, |U^\star|} ,
  \qquad
  r_{\rm term} = \frac{0.2+r}{0.2+r+r^2} , \\
  \mathrm{norm\_term} &=
    \frac{\big\| \Kstar \cdot \grad\driver^\star \big\|}{\norm{\Kstar}_F\,\big\|\grad\driver^\star\big\|}
    \quad\big(\to 1 \text{ if } \grad\driver^\star = 0,\ \text{by convention}\big), \\
  \alpha &= \mathrm{norm\_term}\times r_{\rm term} ,
  \qquad \alpha \to 0 \ \text{ if } U^\star = 0 \text{ or } \norm{\Kstar}_F = 0 .
  \label{eq:alpha_def}
\end{align}
The dimensionless ratio $r$ compares the characteristic HLL numerical
flux scale ($v_{\rm HLL}\,\driver^\star/|\mathbf{dx}_{ij}|$, after
clearing the $|\mathbf{dx}_{ij}|$ factor) against the physical diffusive
flux scale ($\norm{\Kstar}_F\,U^\star/|\mathbf{dx}_{ij}|^2$, similarly
normalised): $r_{\rm term}\to 1$ when $r\gg1$ (numerical transport
dominates: little artificial diffusivity is needed to stabilise it) and
$r_{\rm term}\to 5r$ as $r\to0$ (numerical transport is weak compared to
the physical diffusivity, so the HLL dissipation term is allowed through
more strongly to avoid an under-diffusive, oscillatory update).
\code{norm\_term} additionally attenuates $\alpha$ when $\Kstar$ is
strongly anisotropic relative to the local gradient direction, since
$\|\Kstar\cdot\grad\driver^\star\| \leq \norm{\Kstar}_F
\|\grad\driver^\star\|$ with equality only when $\grad\driver^\star$ lies
along $\Kstar$'s dominant eigendirection.

\begin{definition}[Two-argument minmod]
\label{eq:minmod_def}
\begin{equation}
  \mathrm{minmod}(a,b) =
  \begin{cases}
    \min(a,b) & a>0 \text{ and } b>0 \\
    \max(a,b) & a<0 \text{ and } b<0 \\
    0 & \text{otherwise.}
  \end{cases}
\end{equation}
\end{definition}

\subsection{Direct-Flux Sign-Consistency Guard}
\label{sec:anti-diffusion-guard}
As a final robustness check, the code independently computes a
\emph{direct} (central-difference) estimate of the flux, built purely
from the two particles' own driver values (not their gradients, and not
interface-reconstructed):
\begin{equation}
  q_i = \max(0,\driver_i) , \quad q_j=\max(0,\driver_j) , \quad
  \grad\driver_{\rm dir} = \frac{(q_j-q_i)}{|\mathbf{dx}_{ij,\rm phys}|^2}\,\mathbf{dx}_{ij,\rm phys} ,
\end{equation}
\begin{equation}
  F_{\rm dir} =
  \begin{cases}
    -\big(\Kstar\cdot\grad\driver_{\rm dir}\big)\cdot \dfrac{\mathbf{dx}_{ij,\rm phys}}{|\mathbf{dx}_{ij,\rm phys}|}
      & \text{anisotropic\_gradient mode} \\[1.2em]
    -\bar\kappa\, \grad\driver_{\rm dir}\cdot \dfrac{\mathbf{dx}_{ij,\rm phys}}{|\mathbf{dx}_{ij,\rm phys}|}
      & \text{otherwise, } \bar\kappa=\tfrac12(\kappaP_i+\kappaP_j) .
  \end{cases}
\end{equation}
If this direct estimate disagrees in \emph{sign} with the HLL flux of
Eq.~\eqref{eq:flux_hll} by more than a relative tolerance
$\epsHLL \equiv$ \code{chem\_data->hll\_riemann\_solver\_epsilon},
\begin{equation}
  F_{\rm dir}\cdot F_{\rm HLL} < 0
  \quad\text{and}\quad
  |F_{\rm dir}| > \epsHLL\, |F_{\rm HLL}| ,
\end{equation}
the interface flux is set to zero rather than trusted
(\code{parabolic/chemistry\_riemann\_HLL.h}). This guards against
the minmod-blended HLL construction of Eq.~\eqref{eq:flux_hll}
occasionally returning an up-gradient (anti-diffusive) flux due to
reconstruction error, at the cost of locally stalling diffusion across
that one interface for one sub-step. The final flux entering the update
is $F_{\rm HLL}\times\Aij$ (physical face area,
\code{parabolic/chemistry\_flux.h}).

\section{Flux Limiting for Physical Realizability}
\label{sec:flux-limiter}

Even a Riemann-consistent flux can, once multiplied by a finite timestep
and summed over all of a particle's neighbours, remove more metal mass
than the particle possesses. \code{chemistry\_limit\_metal\_mass\_flux}
(\code{chemistry\_flux.h}) rescales the raw \emph{pairwise} flux
$\Delta m \equiv F_{\rm HLL}\Aij\, \dt_{\min}$ ($\dt_{\min}$ the smaller
of the two particles' active chemistry timesteps) through a sequence of
symmetric (particle order-independent) guards:

\begin{enumerate}
  \item \textbf{Source constraint.} Let $m_{Z}^{\rm src}$ be the current
        metal mass of whichever particle is donating mass (determined
        by $\mathrm{sign}(\Delta m)$, not by particle index; this is
        what makes the limiter symmetric under $i\leftrightarrow j$). If
        $m_Z^{\rm src}\le 0$, the flux is annihilated: a particle cannot
        export metal mass it does not have.
  \item \textbf{Capacity constraint.} Let $m^{\rm sink}$ be the sink
        particle's total (gas) mass and $m_Z^{\rm sink}$ its current
        metal mass; define the remaining headroom before $\Zfrac=1$,
        $C \equiv m^{\rm sink} - m_Z^{\rm sink}$. If $C\le 0$ the flux is
        annihilated.
  \item \textbf{Noise gate.} If
        $|\Delta m| < \epsilon_{\rm noise}\, m_Z^{\rm src}$ with
        $\epsilon_{\rm noise}=10^{-15}$
        (\code{GEAR\_FVPM\_DIFF\_NOISE\_GATE}), the flux is treated as
        floating-point noise and zeroed, preventing a slow
        accumulation ("ratchet effect") of sub-machine-precision
        transfers.
  \item \textbf{Smooth stability/startup limiter.} An effective mass
        budget is formed,
        \begin{equation}
          m_{\rm eff} = \min\!\Big(\min(m_Z^{\rm src}, C),\ \
          \frac{\max\!\big(0.25\, m_Z^{\rm sink},\ 0.1\, m_Z^{\rm src}\big)}{0.5}\Big) ,
        \end{equation}
        combining the source/capacity bound with the larger, via
        $\max(\cdot,\cdot)$, of two further bounds: a
        \emph{sink-stability} bound $0.25\, m_Z^{\rm sink}$
        (\code{GEAR\_FVPM\_DIFF\_LIMITER\_SINK\_STABILITY}: no single
        pairwise interaction may change the sink's metal mass by more
        than $25\%$, preventing oscillatory "neighbour flooding") and a
        \emph{startup} floor $0.1\, m_Z^{\rm src}$
        (\code{GEAR\_FVPM\_DIFF\_LIMITER\_STARTUP}).

        The startup floor exists because the sink-stability bound is
        proportional to $m_Z^{\rm sink}$ itself, so for pristine gas
        ($m_Z^{\rm sink}=0$) that bound is exactly zero. Without the
        $\max$, a zero bound zeroes the flux into that particle on
        every interaction; since a zeroed flux leaves $m_Z^{\rm sink}$
        at zero, this is a fixed point, a particle caught at
        $m_Z^{\rm sink}=0$ could never receive metal at all, stalling
        the diffusion front permanently at the edge of every enriched
        region rather than merely slowing it down. The startup floor
        breaks this deadlock by capping the sink-stability term from
        below at a fraction of the \emph{source's} mass instead,
        independent of how little the sink currently holds.

        Because the two bounds are combined with $\max$, the startup
        floor is the one actually governing the limit exactly while
        $m_Z^{\rm sink} < (0.1/0.25)\,m_Z^{\rm src} = 0.4\,
        m_Z^{\rm src}$, i.e.\ while the sink remains under-enriched
        relative to this particular neighbour; past that ratio the
        ordinary sink-stability bound takes over on its own, with no
        explicit switch. The startup floor is therefore a
        \emph{bootstrap} mechanism, active only at the leading edge of
        an enrichment front, not a persistent scale parameter, and it
        never relaxes the source/capacity bound $\min(m_Z^{\rm src},
        C)$ outside the $\max(\cdot,\cdot)$: however large it is set,
        it cannot enrich a sink beyond what the source actually has
        available. The safety scale $0.5$
        (\code{GEAR\_FVPM\_DIFF\_LIMITER\_SAFETY}) leaves headroom below the
        hard constraints. The flux is then attenuated smoothly rather
        than hard-clipped,
        \begin{equation}
          x = \frac{|\Delta m|}{0.5\, m_{\rm eff} + 10^{-40}} , \qquad
          \Delta m \;\longrightarrow\; \frac{\Delta m}{1+x} .
        \end{equation}
        As $|\Delta m|\to\infty$ this rational form asymptotes to a
        bounded transferred mass $\approx 0.5\,m_{\rm eff}$, while for
        $|\Delta m|\ll 0.5\,m_{\rm eff}$ the attenuation factor is
        $\approx 1$ (unmodified flux), giving a $C^0$-continuous
        transition instead of a hard clip that would otherwise reduce
        the formal order of convergence of the scheme near the
        constraint boundary.
\end{enumerate}

\subsection{Global Flux-Corrected Transport Positivity Rescale}
\label{sec:fct}

The guards above bound each \emph{pairwise} interaction individually, but
a particle can be visited by many active neighbours within the same
timestep. Even when every individual transfer respects its own bound, the
\emph{sum} of many simultaneously limited outflows can still exceed the
donor's available metal mass. A global, per-particle correction closes
this gap, following the flux-corrected transport (FCT) approach of
\citet{Zalesak1979}: a low-order (here, zero) update is guaranteed
positive, and the high-order correction (the diffusive flux itself) is
scaled back just enough to keep the sum non-negative.

\subsubsection{Two-Pass Mechanism}
The correction runs as two passes over the same set of pair interactions,
scheduled around a dedicated ghost task
(\code{runner\_do\_chemistry\_fct\_ghost}, \code{runner\_ghost.c}) that sits
between the density/gradient finalisation and the force loop:
\begin{enumerate}
  \item \textbf{Prep pass.} For every pair interaction, the same
        already limiter-attenuated flux $\Delta m$ of the previous section
        is computed, via a shared, out-of-line function so both passes are
        bit-identical (\code{chemistry\_compute\_pair\_fluxes},
        \code{chemistry.c}), and its donor-side outflow magnitude is
        accumulated per particle and per species,
        \begin{equation}
          \label{eq:fct-sum}
          \Sigma_{\rm out} = \sum_{\rm neighbours} \max(0, -\Delta m) .
        \end{equation}
  \item \textbf{Ghost task.} Once every neighbour's contribution has been
        accumulated, each particle's available metal mass for this step is
        formed from its current mass plus the net flux already committed
        from earlier steps in its active timestep window,
        $m_{\rm avail} = \max(0,\, m_Z + \Delta m_{\rm committed})$, and a
        rescale factor is computed,
        \begin{equation}
          \label{eq:fct-theta}
          \theta = \begin{cases}
            \dfrac{(1-\delta)\, m_{\rm avail}}{\Sigma_{\rm out}} & \Sigma_{\rm out} > (1-\delta)\, m_{\rm avail} \\[0.8em]
            1 & \text{otherwise}
          \end{cases} , \qquad \delta = 10^{-6} .
        \end{equation}
        $\Sigma_{\rm out}$ is then reset to zero for the next step.
  \item \textbf{Apply pass.} The force loop recomputes the identical
        pairwise flux and multiplies it by the frozen, donor-side $\theta$
        before accumulating it into the particle's flux buffer: the final
        flux entering the update is
        \begin{equation}
          \Delta m \;\longrightarrow\; \theta_{\rm donor}\, \Delta m .
        \end{equation}
\end{enumerate}
Because $\theta$ depends only on quantities already known before the force
loop runs, this is a single closed-form rescale, not an iterative
procedure.

\subsubsection{Composition with the Pairwise Limiter}
The two mechanisms are composed, not exclusive. Every pairwise flux is
already attenuated by the smooth limiter above before it ever reaches
Eq.~\eqref{eq:fct-sum}, so a tightly tuned pairwise limiter keeps
$\Sigma_{\rm out}$ well below $m_{\rm avail}$ for most particles,
$\theta\approx 1$ almost everywhere, and the pairwise limiter does
essentially all of the positivity-relevant work. A loosely tuned pairwise
limiter instead lets more flux traffic reach the sum, engaging $\theta<1$
more often. Because $\theta$ is a single per-particle, per-species factor,
engaging it suppresses every outgoing transfer from that donor for that
species that step uniformly, including transfers that individually were
not a problem. The pairwise limiter parameters therefore still control how
much of the positivity burden is carried by this coarser, per-particle
correction versus the pairwise limiter's own targeted attenuation.

\subsubsection{Not an Exact Positivity Proof}
Eq.~\eqref{eq:fct-theta} bounds outflow only: $\Sigma_{\rm out}$ sums the
magnitude of outgoing transfers, while the committed net flux
$\Delta m_{\rm committed}$ sums signed in- and outflow. For a particle
draining toward a small residual mass, the relative round-off of that net
scales with the ratio of total flux traffic to the residual, which is
formally unbounded. The back-off constant $\delta=10^{-6}$ is an empirical
margin: across the module's test suite it accommodates
traffic-to-availability ratios up to about $10^{6}$, typically reducing
the number of particles left with a negative metal mass by three to four
orders of magnitude relative to running with no global correction at all,
but a particle draining to a residual far smaller than that margin can
still be left marginally negative internally.

This residual is not itself the last line of defence. Every physics module
outside the diffusion solver reads metal mass exclusively through
\code{chemistry\_get\_part\_corrected\_metal\_mass}
(Section~\ref{sec:unphysical}), which clamps a negative value to zero on
read. The diffusion solver is therefore free to carry a small transient
negative internally between steps, but no other module, including
cooling, ever observes it.

\textbf{MPI.}
The donor-side sums of Eq.~\eqref{eq:fct-sum} are purely local: a
particle's neighbours across a rank boundary do not currently contribute
to its $\Sigma_{\rm out}$. Rather than silently drop this correction on
multi-rank runs, the code raises a hard error at task-graph construction
time when configured with more than one MPI rank; distributed
accumulation is not yet implemented.

\section{Time Integration}
\label{sec:time-integration}

The scheme is explicit and operator-split from hydrodynamics, gravity,
cooling and star formation/feedback. For each active particle pair, the
signed flux is accumulated into both particles' per-timestep flux
buffers with the sign convention that a positive flux (in the direction
of $\nunit$, from $i$ towards $j$) is subtracted from the left state $i$
and added to the right state $j$
(\code{chemistry\_part\_update\_fluxes\_left/right},
\code{parabolic/chemistry\_flux.h}):
\begin{equation}
  \Delta m_{Z,i} \mathrel{-}= F_{\rm HLL}\,\Aij\, \dt_{\min} , \qquad
  \Delta m_{Z,j} \mathrel{+}= F_{\rm HLL}\,\Aij\, \dt_{\min} .
\end{equation}
Because SWIFT is an individual-timestep code, inactive neighbours of an
active particle must still receive their share of the exchanged flux to
preserve exact pairwise mass conservation; this is guaranteed structurally
by \code{runner\_iact\_chemistry\_fluxes\_common}
(\code{chemistry\_iact.h}), which always updates both sides of a
symmetric interaction. At the end of the force loop, each particle
applies its accumulated buffer to the conserved variable
(\code{chemistry\_end\_force}, \code{chemistry.h}):
\begin{equation}
  m_{Z,i}^{n+1} = m_{Z,i}^{n} + \sum_{j\in\mathcal N(i)} \Delta m_{Z,i}^{(j)} ,
\end{equation}
i.e.\ a first-order (forward Euler) update of the conserved metal mass
using the already-limited, already-integrated-over-$\dt_{\min}$ flux sum:
no further sub-cycling or higher-order time integrator is applied to the
diffusion term itself.

\section{Robustness: Unphysical-State Detection and Correction}
\label{sec:unphysical}

All corrections live in \code{chemistry\_unphysical.h} and are invoked
from exactly three call sites, distinguished by an integer
\code{callloc}, each with a different tolerance for how "wrong" a state
is allowed to be before the code intervenes rather than simply
correcting silently
(\code{chemistry\_check\_unphysical\_state},
\code{chemistry\_unphysical.h}):

\begin{itemize}
  \item \textbf{\code{callloc=0}}, after the density loop
        (\code{chemistry\_end\_density}): a negative metal mass here
        indicates corruption \emph{outside} the diffusion solver (e.g.\
        a bug in feedback/enrichment), so the first occurrence for a
        given particle raises a hard \code{error()}.
  \item \textbf{\code{callloc=1}}, during gradient extrapolation
        (Section~\ref{sec:state-reconstruction}): the correction is
        soft. If the extrapolated mass would go negative, the
        extrapolation is discarded and the pre-extrapolation (0th-order,
        particle-centred) value is used instead; the routine returns a
        \code{char corrected} flag consumed by the caller to also reset
        the reconstructed \emph{state} $U$, not just the mass used for
        the check; the flag is compared with an integer test rather than
        a floating-point comparison, avoiding round-off-triggered false
        negatives. If instead the extrapolated mass exceeds the gas
        mass, it is similarly reverted to the 0th-order value (or, in a
        pathological case, hard-clamped to the gas mass).
  \item \textbf{\code{callloc=2}}, after the force loop
        (\code{chemistry\_end\_force}): a negative mass here means the
        limiters of Sections~\ref{sec:flux-limiter}--\ref{sec:anti-diffusion-guard}
        failed to prevent an overshoot. The particle's per-species
        \code{negativity\_counter} is incremented and a throttled warning
        is issued every $200$ consecutive occurrences
        (\code{GEAR\_FVPM\_DIFF\_NEGATIVITY\_COUNTER\_PRINT\_LIMIT}); if the
        magnitude of the negative mass ever exceeds the particle's gas
        mass, a hard \code{error()} is raised regardless of the
        counter, since that indicates a large, likely unrecoverable,
        conservation failure rather than round-off.
\end{itemize}

In addition to this three-site dispatcher:
\begin{itemize}
  \item \textbf{Final safety clamp} (Section~\ref{sec:state-reconstruction},
        step 6): an unconditional
        $U \gets \max(0,U)$ applied to both reconstructed states
        immediately before they are handed to the Riemann solver,
        independent of whether the softer \code{callloc=1} logic fired.
  \item \code{chemistry\_check\_unphysical\_metallicity}
        (\code{chemistry\_unphysical.h}) hard-clamps any
        reconstructed metallicity to $[0,1]$, used when the diffusion
        driver is $\Zfrac$.
  \item \code{chemistry\_check\_unphysical\_total\_metal\_mass}
        (\code{chemistry\_unphysical.h}) checks, after both the
        density loop and the force loop, that the \emph{sum} over all
        tracked species does not exceed the particle's gas mass; if it
        does, every species is rescaled by $1/(10\times m_{Z,\rm
        tot}/m)$ and a warning is issued (this is a diagnostic
        safety net, not expected to trigger in a correctly configured
        run).
  \item \code{chemistry\_get\_part\_corrected\_metal\_mass}
        (\code{chemistry\_getters.h}) is the \emph{only} interface
        exposed to other physics modules (feedback, cooling, star
        formation): it always clamps negative metal mass to zero and
        caps $m_Z \geq m$ down to $0.01\,m$. The diffusion solver is
        thus free to carry small transient negative or super-unity
        states internally (a deliberate design choice to avoid
        amplifying noise, see the comments in
        Section~\ref{sec:gradients}), while guaranteeing that no other
        module ever observes an unphysical metallicity.
\end{itemize}

\section{Timestep Stability Criteria}
\label{sec:timestep}

\subsection{Characteristic Particle Size: GIZMO and Lanson \& Vila Definitions}
\label{sec:particle-size}
The characteristic length entering every CFL estimate is a particle's
own resolution scale $\dx$, computed one of two ways depending on a
compile-time switch; both describe the same physical scale, so the
choice is one of numerical convenience, not physics (see
\textbf{Equivalence} below).

\textbf{GIZMO definition (default).} The physical GIZMO particle size,
\begin{equation}
  \dx = a \left(\frac{V_i}{\Omega_d}\right)^{1/d} ,
\end{equation}
with $\Omega_d$ the volume of the $d$-dimensional unit sphere
(\code{chemistry\_get\_physical\_particle\_size}, \code{chemistry\_timesteps.h}).
This is the same getter the hyperbolic solver's causal bound uses, so
the CFL length and the causal bound's length scale cannot drift apart.

\textbf{Lanson \& Vila definition (alternative).} Under
\code{GIZMO\_LANSON\_VILA\_PARTICLE\_SIZE},
\begin {equation}
  \dx = a\,\delta x_i \qquad \delta x_i = \frac{1}{2\sum_l w_l\|A_{il}\|}
\end{equation}
is used instead, with $\delta x_i$ built from the same face-area
accumulator as Eq.~\eqref{eq:face_area}. Their eq.~(58) defines
$\delta x_i$ with $w_l$ the neighbour's \emph{volume}, evaluated here
as $\sum_l\|A_{il}\|/V_l$. Notice that the kernel
value $w_jh_j^{-d}$ is dimensionally interchangeable with $1/V_l$ but
smaller by a factor of order the neighbour count, and using it in that
place would inflate $\delta x_i$ by the same factor.

\textbf{Equivalence.} The authors note below eq.~(58) that
$\sum_l w_l\|A_{il}\| = O(1/\epsilon_0)$, with $\epsilon_0$ the
smoothing length, so the two definitions describe the same resolution
scale and differ only by an $O(1)$ constant.

\textbf{Consequence for $\CCFL$.} Switching between them at a fixed
$\CCFL$ therefore buys nothing by itself: the constant is what a
Courant number absorbs, so the admissible $\CCFL$ has to be re-derived
for whichever length is in use, from Lanson \& Vila's own
operator-based stability condition (their eq.~62) rather than carried
over.

\subsection{Mode-Dependent CFL Conditions}
\label{sec:cfl-modes}
The three diffusion modes do \emph{not} share a single timestep formula
(\code{chemistry\_diffusion\_timestep}, \code{chemistry\_timesteps.h}).
In every mode the Frobenius norm $\norm{\Kmat}_F$ is evaluated first; if
it vanishes (no diffusivity yet computed, e.g.\ at start-up) the
timestep constraint is skipped ($\dt\to\infty$).

\paragraph{Isotropic constant.}
\begin{equation}
  \label{eq:dt_const}
  \dt \le \CCFL\, \frac{\dx^2}{\norm{\Kmat}_F} = \CCFL\, \frac{\dx^2}{\kappaC} ,
\end{equation}
the textbook explicit-diffusion CFL bound (derived rigorously in
Appendix~\ref{app:von-neumann}), with no additional density-dependent
prefactor since $\driver=\rhoZ$ already carries the density.

\paragraph{Isotropic Smagorinsky and anisotropic gradient.}
Both go through the more elaborate form (Hopkins~\citeyear{Hopkins2017diffusion}, eq.~15-style)
\begin{equation}
  \label{eq:dt_aniso}
  \dt \le \CCFL\, \frac{E^2}{\norm{\Kmat}_F}\,\bar\rho ,
\end{equation}
where the extra factor $\bar\rho$ arises because here the conserved
state is $U=\rhoZ$ but the driver is $\driver=\Zfrac$, so
$\norm{U}/\norm{\driver} = \bar\rho$ (physical gas density,
\code{hydro\_get\_physical\_density}), and $E$ is a gradient-strength-aware
effective length,
\begin{equation}
  \label{eq:E-bracket}
  E =
  \begin{cases}
    \dx & \text{isotropic\_smagorinsky (see remark below)} \\[4pt]
    \dfrac{\norm{\driver}\,\dx}{\norm{\grad\driver}\,\dx + \norm{\driver}}
      & \text{anisotropic\_gradient,}
  \end{cases}
\end{equation}
with $\norm{\driver} = \big(\sum_{\rm species} \Zfrac^2\big)^{1/2}$ and
$\norm{\grad\driver}$ the Frobenius norm of the stacked
species-by-component metallicity gradient matrix, both evaluated
particle-locally. If $\norm{\driver}=0$ (no metal to diffuse) the
constraint is skipped.

For the anisotropic mode, $E$ interpolates between two regimes: in a
smooth region ($\norm{\grad\driver}\,\dx \ll \norm{\driver}$),
$E \to \dx$ and Eq.~\eqref{eq:dt_aniso} reduces to the same
$\dx^2/\norm{\Kmat}_F$ scaling as Eq.~\eqref{eq:dt_const}. In a region
with a locally steep metallicity gradient
($\norm{\grad\driver}\,\dx \gg \norm{\driver}$),
$E \to \norm{\driver}/\norm{\grad\driver}$, the gradient's own local
scale length, which can be much shorter than $\dx$; the timestep
constraint then becomes correspondingly \emph{tighter} than the naive
particle-size estimate. This is intentional and physically necessary:
diffusion locally erases steep features fastest exactly where they are
steep, but the code comment at \code{chemistry\_timesteps.h}
notes explicitly that it is expensive in practice: strong feedback-driven
enrichment gradients can drive this term to very small values and
"slow down everything." For that reason the Smagorinsky mode
deliberately \emph{does not} use the bracket of
Eq.~\eqref{eq:E-bracket}, substituting the plain particle size $\dx$
instead (\code{chemistry\_timesteps.h}), a pragmatic,
non-rigorous compromise flagged as such directly in the source, adopted
because the physically correct bound was found to be prohibitively
restrictive for that closure in practice.

\section{Validation and Verification}

\subsection{Conservation Properties}
The pairwise sign convention of Section~\ref{sec:time-integration}
(flux subtracted from the source, added to the sink, both particles
updated from a single Riemann solve) guarantees exact discrete
conservation of total metal mass across every interacting pair,
independent of the flux-limiter stage of Section~\ref{sec:flux-limiter}
(the limiter rescales the flux by a common factor \emph{before} it is
split between the two particles, so the pairwise anti-symmetry is
preserved). Global conservation can therefore only be broken by the
\code{callloc=2} corrections of Section~\ref{sec:unphysical}, which are
intentionally non-conservative safety nets and are designed to trigger
rarely (tracked via the negativity counter for exactly this reason).

\subsection{Stability}
Section~\ref{sec:timestep}'s CFL bounds are necessary but, because of
the explicit minmod/artificial-diffusivity blending in
Eq.~\eqref{eq:flux_hll} and the nonlinear flux limiter of
Section~\ref{sec:flux-limiter}, not proven sufficient in closed form for
the full nonlinear scheme; Appendix~\ref{app:von-neumann} derives the
linear (constant-coefficient, unlimited) stability bound that motivates
Eq.~\eqref{eq:dt_const}.

\subsection{Suggested Test Problems}
\begin{itemize}
  \item \textbf{1D/3D diffusion of an isolated metal peak}
        (\code{isotropic\_constant} mode): compare the discrete solution
        against the analytic Green's function of the linear heat
        equation, Eq.~\eqref{eq:pde_const}, on a uniform particle
        distribution.
  \item \textbf{Two-peak merging test}: verify that two well-separated
        enriched regions diffuse independently at early times and merge
        smoothly, exercising the cell and face limiters of
        Section~\ref{sec:gradients} in a regime with two simultaneous
        local extrema.
  \item \textbf{Shear-flow anisotropic diffusion}: impose a prescribed
        background shear (e.g.\ a Couette or Kolmogorov flow) and verify
        that metal transport is preferentially aligned with the
        positive-eigenvalue directions of $\Splus$
        (Proposition~\ref{prop:psd}), not merely isotropic.
\end{itemize}


% theory/ParabolicDiffusion/hyperbolic.tex
% Body of the hyperbolic (Cattaneo-relaxation) FVPM diffusion scheme.
% Input from GEAR_FVPM_Diffusion_Theory.tex (the master document); not
% standalone.

\part{Hyperbolic (Cattaneo-Relaxation) Diffusion}
\label{part:hyperbolic}

\section{Formulation}

\subsection{Purpose and Scope}
\label{sec:scope}

The parabolic scheme of Part~\ref{part:parabolic} links the diffusive
flux algebraically to the local gradient (Fick's law,
Eq.~\eqref{eq:ficks_law}): any change in $\grad\driver$ is felt
instantaneously and everywhere, an infinite propagation speed with no
physical counterpart in real turbulent mixing. It is also numerically
awkward precisely where the gradient is steepest and least resolved,
e.g.\ at a supernova injection site, where one timestep's worth of freshly
enriched gas sits next to pristine neighbours: the parabolic CFL bound
tightens without limit as the gradient steepens
(Section~\ref{sec:speed-criterion-aniso}). The hyperbolic
(Cattaneo~\citeyear{Cattaneo1948}-relaxation) scheme, selected at compile
time by \code{CHEMISTRY\_GEAR\_FVPM\_HYPERBOLIC\_DIFFUSION}
(\code{src/chemistry/GEAR\_FVPM\_DIFFUSION/hyperbolic/}), fixes both
problems by promoting the flux itself to a dynamical variable that
relaxes toward its Fickian value on a finite timescale $\taurelax$
(Section~\ref{sec:governing-system}), giving the system a genuine finite
characteristic speed $\chyp$ and an ordinary wave-speed CFL condition in
place of the parabolic diffusion-time bound.

This part starts with a first-principles derivation of that
characteristic speed, promised but not derived where
Section~\ref{sec:parabolic-scope} of Part~\ref{part:parabolic} first
mentions the system, and a direct confrontation of that result against
the two closed-form expressions currently returned by
\code{chemistry\_get\_physical\_hyperbolic\_soundspeed}
(\code{hyperbolic/chemistry\_getters.h}). It then works outward to the
practical consequences: how the relaxation-time normalisation
$\taurelax_0$ must be chosen relative to the diffusion coefficient
(Section~\ref{sec:tau0-physical}), how an overestimated or unbounded wave
speed leaks into excess numerical dissipation and how that is controlled
(Sections~\ref{sec:causal-bound}--\ref{sec:hyperbolic-limiter}), how the
stiff relaxation source term is integrated in the half-timestep predictor
(Section~\ref{sec:predictor}), and finally the exact Green's functions
the scheme can be validated against (Section~\ref{sec:exact-solutions}).

\subsubsection{Relation to the Parabolic Part}
This part assumes familiarity with the FVPM geometry, the diffusion
driver $\driver$, and the diffusion tensor $\Kmat$ as defined in
Part~\ref{part:parabolic}, Sections~\ref{sec:physical-foundations}--
\ref{sec:fvpm-geometry}; it does not repeat that material. The two parts
share notation (see the macro list in the preamble) so that symbols mean
the same thing in both, and they share the FVPM geometry, gradient
estimator and slope limiters of Part~\ref{part:parabolic},
Sections~\ref{sec:fvpm-geometry}--\ref{sec:gradients}. They differ in
everything downstream of the reconstructed states: different Riemann
solvers, different conserved-state vectors ($U_0=\rhoZ$ alone for the
parabolic scheme versus the four-component $\Uvec=(\rhoZ,\Fdiff)$ here),
and different timestep criteria (Section~\ref{sec:speed-criterion}, a
genuine wave-speed CFL condition, rather than
Part~\ref{part:parabolic}, Section~\ref{sec:cfl-modes}'s parabolic
diffusion-time bound).

\subsection{Governing System}
\label{sec:governing-system}

\subsubsection{Conservation of Metal Mass}
As in the parabolic case (Part~\ref{part:parabolic},
Eq.~\eqref{eq:parabolic_diffusion_rho_Z}), conservation of metal-mass
density in the frame comoving with each Lagrangian particle reads
\begin{equation}
  \label{eq:mass_conservation}
  \pdv{(\rhoZ)}{t} + \div \Fdiff = 0 .
\end{equation}
The difference between the parabolic and hyperbolic schemes is entirely in
what closes this equation: the parabolic scheme links $\Fdiff$
algebraically to $\grad\driver$ via Fick's law
(Eq.~\eqref{eq:ficks_law}); the hyperbolic scheme instead promotes
$\Fdiff$ to an independent dynamical variable, evolved by its own PDE.

\subsubsection{The Cattaneo Relaxation Equation}
The flux obeys:
\begin{equation}
  \label{eq:cattaneo}
  \pdv{\Fdiff}{t} = - \frac{\Kmat}{\taurelax}\grad\driver - \frac{\Fdiff}{\taurelax} ,
\end{equation}
where $\taurelax = \code{p->chemistry\_data.tau}$ is the local relaxation
time (Section~\ref{sec:tau-modes}) and $\driver$ is the same diffusion
driver as in the parabolic scheme ($\driver=\rhoZ$ for
\code{isotropic\_constant}; $\driver=\Zfrac$ otherwise: see
\code{chemistry\_compute\_flux}, \code{hyperbolic/chemistry\_flux.h}).
As $\taurelax\to0$, Eq.~\eqref{eq:cattaneo} tends to
$\Fdiff\to-\Kmat\grad\driver$, i.e.\ it collapses onto Fick's law and
Eq.~\eqref{eq:mass_conservation} becomes the parabolic equation.

The pair \eqref{eq:mass_conservation}--\eqref{eq:cattaneo} is exactly the
Cattaneo/Maxwell-Cattaneo regularisation of a diffusion law, structurally
identical to the M1 closure of radiative transfer
(Berthon et al.~\citeyear{Berthon2007}).

\begin{remark}[$\grad\driver$ is always evaluated at the particle, never extrapolated]
\label{rem:cell-centred-gradients-only}
Every use of $\Kmat\grad\driver$ in this scheme (Eq.~\eqref{eq:cattaneo}
itself, and $\Fdiff_{\rm parabolic}=-\Kmat\grad\driver$, the relaxation
target inside \code{chemistry\_part\_integrate\_flux\_source\_term}) uses
$\grad\driver$ as stored at the particle's own position, never
extrapolated to an interface. This is a choice of implementation to reduce computation cost. Extrapolating
$\grad\driver$ itself to an interface (as opposed to extrapolating
$\driver$ using the already-computed $\grad\driver$, which \emph{is} done
during state reconstruction) would require the gradient of
$\grad\driver$, i.e.\ a second derivative that would be required for each metal species.
Cell-centred evaluation of $\Fdiff_{\rm parabolic}$ is therefore the only
option available with a first-order (linear) gradient estimator, not an
approximation chosen among alternatives.
\end{remark}

\subsubsection{Relaxation-Time Models}
\label{sec:tau-modes}
Two independent modes exist
(\code{chemistry\_compute\_physical\_tau}, \code{hyperbolic/chemistry\_getters.h}):
\begin{align}
  \text{\code{constant\_mode}:} \qquad & \taurelax = \taurelax_0 , \\
  \text{\code{turbulent\_mode}:} \qquad & \taurelax = \frac{\taurelax_0}{\kappaC\,\norm{\Shear}_F} ,
  \label{eq:tau_turbulent}
\end{align}
with $\taurelax_0=$\code{chem\_data->tau} the parameter-file normalisation
and $\Shear$ the (unregularised) physical velocity shear tensor. The
\code{turbulent\_mode} formula includes an explicit $\kappaC^{-1}$
(\code{chem\_data->diffusion\_coefficient}) coupling, motivated and derived
in Section~\ref{sec:tau0-physical}. This choice of relaxation-time model is
independent to the diffusion-tensor mode $\Kmat$
(Section~\ref{sec:diffusion-tensor-recap}): \code{diffusion\_mode}.

\subsubsection{Diffusion Tensor Modes (recap)}
\label{sec:diffusion-tensor-recap}
From \code{chemistry\_get\_physical\_matrix\_K}/\code{chemistry\_get\_diffusion\_coefficient}
(\code{chemistry\_getters.h}):
\begin{align}
  \text{\code{isotropic\_constant}:} \qquad & \Kmat=\kappaP\mathbf I, \quad \kappaP=\kappaC, \\
  \text{\code{isotropic\_smagorinsky}:} \qquad & \Kmat=\kappaP\mathbf I, \quad \kappaP=\kappaC\,\gammak^2\,h^2\,\rho\,\norm{\Shear}_F, \\
  \text{\code{anisotropic\_gradient}:} \qquad & \Kmat = \kappaC\,\gammak^2\,h^2\,\rho \; \Splus .
\end{align}
Crucially, $\Kmat$ carries an explicit factor of $\rho$ in the
\code{isotropic\_smagorinsky} and \code{anisotropic\_gradient} modes, but
\emph{not} in \code{isotropic\_constant}. This asymmetry is the origin of
the consistency check of Section~\ref{sec:confrontation}.

\subsubsection{The Effective Diffusivity $\Deff$}
\label{sec:effective-diffusivity}
$\Kmat$ is defined with respect to the \emph{driver} $\driver$ (Fick's
law, $\Fdiff=-\Kmat\grad\driver$), not with respect to the
\emph{conserved} state $U_0=\rhoZ$, and, as Section~\ref{sec:tau-modes}
already had to confront for $\taurelax$, $\driver$ itself is
\code{diffusion\_mode}-dependent ($\driver=U_0$ for
\code{isotropic\_constant}, $\driver=U_0/\rho$ otherwise). The \code{"DiffusionMatrices"}
snapshot field's own description (\code{chemistry\_io.h}) states it
outright: \emph{``The effective diffusivity is defined as $\Deff =
\Kmat\driver/U$.''} Formally, under the frozen-coefficient approximation
of Section~\ref{sec:quasilinear} ($\driver=\mu U_0$, $\mu$ locally
constant),
\begin{equation}
  \label{eq:D_definition}
  \Deff \equiv \mu\,\Kmat = \frac{\driver}{U_0}\,\Kmat ,
  \qquad
  \mu = \begin{cases} 1 & \code{isotropic\_constant} \\ 1/\rho & \text{otherwise} \end{cases} .
\end{equation}
This is exactly the same $q/U$ (or $U/q$) bookkeeping that recurs
throughout Hopkins~\citeyear{Hopkins2017diffusion} and is already used
elsewhere in this module for precisely this reason. $\Deff$, not $\Kmat$, is the tensor that
plays the role of an ordinary Fickian diffusivity \emph{for the conserved
variable $U_0$ directly}: substituting $\grad\driver\approx\mu\grad U_0$
(frozen-coefficient) into Fick's law gives $\Fdiff=-\Kmat\mu\grad U_0 =
-\Deff\grad U_0$, i.e.\ $\Deff$ is what $\Kmat$ becomes once rewritten
entirely in terms of $U_0$ instead of $\driver$. 

\section{Characteristic Structure}

\subsection{Quasi-Linear Form and the Flux Jacobian}
\label{sec:quasilinear}

\subsubsection{State Vector and Frozen Coefficients}
Define the four-component hyperbolic state vector
\begin{equation}
  \Uvec = (U_0, U_1, U_2, U_3) \equiv (\rhoZ, F_x, F_y, F_z) .
\end{equation}
As is standard for deriving a local wave-speed estimate for a Godunov-type
solver, we treat $\Kmat$ (equivalently $\Deff$), $\taurelax$ and $\rho$
as \emph{locally frozen} coefficients for the purpose of constructing the
flux Jacobian: they are evaluated at the particle/interface, not
differentiated with respect to $\Uvec$, matching how $\chyp$ is itself
evaluated from each particle's own, locally stored $\Kmat,\taurelax,\rho$
(\code{chemistry\_iact.h}). Under this approximation,
$\driver=\mu U_0$ (Section~\ref{sec:effective-diffusivity})
is linear in $U_0$, and $\Deff=\mu\Kmat$ is itself frozen.

\subsubsection{The Flux Tensor, Rewritten in Terms of $\Deff$}
The flux of $\Uvec$ in Cartesian direction $k\in\{x,y,z\}$, computed by
\code{chemistry\_get\_hyperbolic\_flux}
(\code{hyperbolic/chemistry\_getters.h}), is $(F_k,\,K_{xk}\driver/\taurelax,\,
K_{yk}\driver/\taurelax,\,K_{zk}\driver/\taurelax)$. Substituting
$\driver=\mu U_0$ (frozen $\mu$) turns
$K_{jk}\driver/\taurelax$ into $D_{jk}U_0/\taurelax$, so
\begin{equation}
  \Gflux{k}(\Uvec) = \left( F_k, \; \frac{D_{xk}U_0}{\taurelax}, \; \frac{D_{yk}U_0}{\taurelax}, \; \frac{D_{zk}U_0}{\taurelax} \right) ,
\end{equation}
i.e.\ the flux tensor is entirely in terms of $\Deff$ and $U_0$, the two
quantities the state vector $\Uvec$ and its own conservation law are
already written in terms of: no leftover $\driver$/$\mu$ bookkeeping.
Eqs.~\eqref{eq:mass_conservation}--\eqref{eq:cattaneo} read, in
conservative form, $\partial_t\Uvec + \sum_k \partial_{x_k}\Gflux{k}(\Uvec)
= -(0,\Fdiff/\taurelax)$.

\subsubsection{The Directional Flux Jacobian}
For a Riemann problem at an interface with unit normal $\nunit$, the
relevant $4\times4$ matrix is the projected flux Jacobian
$\mathbf A(\nunit) \equiv \partial(\Gflux{\cdot}(\Uvec)\cdot\nunit)/\partial\Uvec$.
Writing $\mathbf v \equiv \Deff\nunit/\taurelax$ (a 3-vector), direct
differentiation of the flux tensor above gives, for any $\mathbf x =
(x_0,x_1,x_2,x_3)$,
\begin{equation}
  \label{eq:jacobian_action}
  \big(\mathbf A(\nunit)\,\mathbf x\big)_0 = \nunit\cdot(x_1,x_2,x_3) ,
  \qquad
  \big(\mathbf A(\nunit)\,\mathbf x\big)_j = v_j\, x_0 \quad (j=1,2,3) .
\end{equation}

\subsection{The Characteristic Speed: General Result}
\label{sec:main-theorem}

\begin{proposition}[Characteristic speeds of the hyperbolic diffusion system]
\label{prop:main}
For any symmetric positive semi-definite $3\times3$ effective diffusivity
$\Deff$ (Section~\ref{sec:effective-diffusivity}; no assumption on its
shape: isotropic, Smagorinsky-scaled, or fully anisotropic, and no
assumption on which \code{diffusion\_mode} produced it), any
relaxation-time model, and any unit interface normal $\nunit$, the
projected flux Jacobian $\mathbf A(\nunit)$ of Eq.~\eqref{eq:jacobian_action}
has eigenvalues
\begin{equation}
  \label{eq:main_result}
  \lambda \in \left\{ \; +\chyp(\nunit),\; -\chyp(\nunit),\; 0,\; 0 \; \right\} ,
  \qquad
  \boxed{\;\chyp(\nunit) = \sqrt{\dfrac{\nunit^{\!\top}\Deff\,\nunit}{\taurelax}}\;} .
\end{equation}
\end{proposition}

\begin{proof}
Let $\lambda\neq0$ be an eigenvalue with eigenvector $\mathbf x$. From the
second block of Eq.~\eqref{eq:jacobian_action}, $\lambda x_j = v_j x_0$ for
$j=1,2,3$, i.e.\ $x_j = (v_j/\lambda)\,x_0$ (this requires $x_0\neq0$: if
$x_0=0$ then $x_j=0$ for all $j$ from this relation, i.e.\ $\mathbf
x=\mathbf 0$, not a valid eigenvector). Substituting into the first block,
$\lambda x_0 = \nunit\cdot(x_1,x_2,x_3) = (\nunit\cdot\mathbf v/\lambda)\,x_0$,
and since $x_0\neq0$,
\begin{equation}
  \lambda^2 = \nunit\cdot\mathbf v = \frac{\nunit\cdot(\Deff\nunit)}{\taurelax} = \frac{\nunit^{\!\top}\Deff\nunit}{\taurelax} \; \geq 0
\end{equation}
(non-negative because $\Deff=\mu\Kmat$ is PSD: $\mu>0$ and $\Kmat$ PSD),
giving the two roots
$\lambda=\pm\chyp(\nunit)$ of Eq.~\eqref{eq:main_result}. The remaining two
eigenvalues are $\lambda=0$: any $\mathbf x=(0,x_1,x_2,x_3)$ with
$\nunit\cdot(x_1,x_2,x_3)=0$ (a two-dimensional subspace, the plane
transverse to $\nunit$) satisfies $\mathbf A(\nunit)\mathbf x=\mathbf 0$
by direct substitution into Eq.~\eqref{eq:jacobian_action}.
\end{proof}

\begin{remark}
The two zero eigenvalues correspond to the flux components transverse to
the interface normal: they are advected with zero characteristic speed in
this frozen-coefficient, one-dimensional-projection sense (they do not
"propagate" through the Riemann fan). This is the same structure as the
transverse magnetic field components in a 1D MHD Riemann problem, or the
transverse momentum components in a 1D gas-dynamics Riemann problem: real,
but not part of the acoustic-family wave-speed bound. The code's HLL
solver (Section~\ref{sec:confrontation}) uses the same $\lambda_\pm$ bound
for all four flux components (\code{chemistry\_riemann\_HLL.h}),
which is the standard, conservative simplification, not a further
approximation beyond what Proposition~\ref{prop:main} already licenses for
the two genuinely propagating families.
\end{remark}

Proposition~\ref{prop:main} was derived using only
Eqs.~\eqref{eq:mass_conservation}--\eqref{eq:cattaneo} and the
frozen-coefficient flux Jacobian: no assumption was made anywhere about
the internal structure of $\Kmat$ (isotropic, Smagorinsky, or
anisotropic). The two corollaries below specialise it.

\subsubsection{Corollary: Isotropic Recovery}
\begin{corollary}
\label{cor:isotropic}
$\Kmat=\kappaP\mathbf I$ in both \code{isotropic\_constant} and
\code{isotropic\_smagorinsky} modes, but the driver $\driver$ that fixes
$\kappaP$'s \emph{units} (and hence $\mu=\driver/U_0$) differs between
them (Section~\ref{sec:quasilinear}):
\begin{align}
  \driver &= U_0 \quad (\mu=1)\quad \text{for \code{isotropic\_constant}}
  && \big[\kappaP\ \text{has units}\ \text{length}^2\,\text{time}^{-1}\big] , \\
  \driver &= U_0/\rho \quad (\mu=1/\rho)\quad \text{for \code{isotropic\_smagorinsky}}
  && \big[\kappaP\ \text{has units}\ \text{mass}\,\text{length}^{-1}\text{time}^{-1}\big] .
\end{align}
Consequently $\Deff=\mu\kappaP\mathbf I$ is isotropic in both modes,
$\nunit^{\!\top}\Deff\nunit = \mu\kappaP$ for \emph{every} unit vector
$\nunit$ (direction-independence is universal, guaranteed by
Proposition~\ref{prop:main} for any isotropic $\Kmat$ regardless of
$\mu$), and
\begin{equation}
  \chyp = \sqrt{\frac{\mu\,\kappaP}{\taurelax}} ,
  \qquad
  \mu = \begin{cases} 1 & \text{\code{isotropic\_constant}} \\
                       1/\rho & \text{\code{isotropic\_smagorinsky}} \end{cases} ,
\end{equation}
\emph{not} independent of which closure $\kappaP$ came from: the two
modes' $\Deff$ differ by exactly the factor $\mu$, even though the
underlying $\Kmat=\kappaP\mathbf I$ formula looks superficially similar.
\end{corollary}
The units mismatch is not a detail specific to $\chyp$: any code path
that consumes $\Kmat$ directly (rather than $\Deff$) must know which
$\driver$ convention is in force for the active \code{diffusion\_mode}, or
it silently mixes an ordinary diffusivity with a mass diffusivity.

\subsubsection{Corollary: The Frobenius Norm Is a (Non-Tight) Upper Bound}
\begin{corollary}
\label{cor:frobenius}
For any unit $\nunit$ and any symmetric PSD $\Kmat$ with eigenvalues
$\lambda_1\geq\lambda_2\geq\lambda_3\geq0$,
\begin{equation}
  \nunit^{\!\top}\Kmat\nunit \; \leq \; \lambda_1 \; \leq \; \norm{\Kmat}_F ,
  \qquad
  \frac{\norm{\Kmat}_F}{\lambda_1} \in [1,\sqrt3\,] ,
\end{equation}
with the upper end of the ratio, $\sqrt3$, attained \emph{only} in the
isotropic case $\lambda_1=\lambda_2=\lambda_3$, and the lower end, $1$,
attained only for a rank-1 (maximally anisotropic) $\Kmat$.
\end{corollary}
\begin{proof}
Write $\nunit=\sum_i c_i\mathbf e_i$ in the eigenbasis of $\Kmat$
($\sum_ic_i^2=1$); then $\nunit^{\!\top}\Kmat\nunit=\sum_i\lambda_ic_i^2
\leq\lambda_1\sum_ic_i^2=\lambda_1$, with equality iff $\nunit=\mathbf
e_1$. Also $\norm{\Kmat}_F^2=\sum_i\lambda_i^2\geq\lambda_1^2$, with
equality iff $\lambda_2=\lambda_3=0$ (rank 1), so $\norm{\Kmat}_F/\lambda_1
\geq1$. For the upper end, in 3D $\sum_i\lambda_i^2 \leq
3\max_i\lambda_i^2=3\lambda_1^2$ with equality iff all three eigenvalues
are equal, giving $\norm{\Kmat}_F/\lambda_1\leq\sqrt3$.
\end{proof}
\begin{example}[Why the Frobenius norm overestimates, in numbers]
\label{ex:frobenius-numbers}
The Frobenius norm $\norm{\Kmat}_F=\sqrt{\sum_{ij}K_{ij}^2}$ is a
different mathematical object from ``the diffusion strength along a
specific direction $\nunit$'', $\nunit^{\!\top}\Kmat\nunit$: it sums
contributions from \emph{every} entry of the matrix, not just the ones
relevant to one direction. Two cases make the gap concrete:
\begin{itemize}
\item \textbf{Isotropic}, $\Kmat=\mathrm{diag}(1,1,1)$ (eigenvalues
  $1,1,1$): the true wave speed in \emph{any} direction uses
  $\nunit^{\!\top}\Kmat\nunit=1$: e.g.\ $\nunit=(1,0,0)$ gives
  $1\cdot1\cdot1+0+0=1$. But $\norm{\Kmat}_F=\sqrt{1^2+1^2+1^2}=\sqrt3\approx1.73$.
  The code would have used $\chyp\propto\sqrt{1.73}$ instead of
  $\chyp\propto\sqrt1$, a $\sim32\%$ overestimate of the speed itself,
  and (Section~\ref{sec:hll-consequence}) the same overestimate directly
  in the HLL numerical dissipation. $\lambda_1(\Kmat)=1$ here: exact, no
  overestimate at all.
\item \textbf{Anisotropic}, $\Kmat=\mathrm{diag}(4,1,1)$ (one strong
  direction, two weak ones): along the strong axis $\nunit=(1,0,0)$, the
  true value is $\nunit^{\!\top}\Kmat\nunit=4$. Here
  $\norm{\Kmat}_F=\sqrt{16+1+1}=\sqrt{18}\approx4.24$, only a
  $\sim6\%$ overestimate, much smaller than the isotropic case, exactly
  as Corollary~\ref{cor:frobenius} predicts ($\norm{\Kmat}_F/\lambda_1\to1$
  as $\Kmat$ becomes more anisotropic). $\lambda_1(\Kmat)=4$: exact again,
  for this direction (and a safe, tighter-than-Frobenius bound for any
  other direction a neighbour might be in).
\end{itemize}
The pattern is the opposite of what intuition might suggest: the
Frobenius norm is worst (loosest) for isotropic $\Kmat$, where directional
information matters least, and best (tightest, in fact exact) for
rank-1 $\Kmat$, where it matters most. $\lambda_1(\Kmat)$ is exact in the
isotropic case and a safe, tighter bound in the anisotropic case; it
dominates $\norm{\Kmat}_F$ everywhere.
\end{example}
Using the Frobenius norm in place of the direction-projected
$\nunit^{\!\top}\Kmat\nunit$ overestimates the true characteristic speed
by a factor of exactly $\sqrt3\approx1.73$ in the isotropic limit
(Example~\ref{ex:frobenius-numbers}), and by a smaller (but never
negative) factor for anisotropic $\Kmat$. It is always a \emph{safe}
(CFL-stable) overestimate, never an underestimate, but it is also never
tight except for rank-1 $\Kmat$. \textbf{The correction is
$\lambda_1(\Kmat)$}, the matrix's largest eigenvalue (Remark~\ref{rem:lambda-max-not-n}
below): swap \code{chemistry\_get\_matrix\_norm(K)} (Frobenius) for
\code{chemistry\_get\_matrix\_max\_eigenvalue(K)} wherever $\chyp$ is
computed.

\begin{remark}[$\lambda_1(\Kmat)$, not $\nunit$, is the right practical fix]
\label{rem:lambda-max-not-n}
Proposition~\ref{prop:main} is stated per-interface, with $\chyp(\nunit)$
depending on the pair-specific normal $\nunit$. It is tempting to conclude
that fixing $\chyp$ in the code requires threading $\nunit$ into
\code{chemistry\_get\_physical\_hyperbolic\_soundspeed} and evaluating it
per-interaction rather than once per particle, but this is unnecessary,
and would be a worse engineering choice. What the CFL condition and the
HLL wave-speed bound actually need is not the exact directional value, but
\emph{any} bound that is never smaller than $\nunit^{\!\top}\Kmat\nunit$
for every $\nunit$ a particle might face. Corollary~\ref{cor:frobenius}
already proves $\lambda_1(\Kmat)$ (the largest eigenvalue of a
particle's own $\Kmat$) is exactly such a bound, and a strictly tighter
one than $\norm{\Kmat}_F$ (equal to it only for isotropic $\Kmat$; strictly
smaller, and exact for the worst-case direction, otherwise). $\lambda_1(\Kmat)$
depends only on $\Kmat$, hence only on particle-local quantities ($\rho$,
$h$, $\Shear$ or $\Splus$): it \emph{is} a particle property, computable
once per particle from an eigendecomposition of $\Kmat$ (SWIFT already has
one, \code{sym\_matrix\_diagonalise\_3x3\_d}, used elsewhere in this module
by \code{chemistry\_regularize\_shear\_tensor}), exactly like $\taurelax$,
$\kappaP$, or the old $\norm{\Kmat}_F$ it replaces. No call site needs a
new argument. The only thing given up relative to the fully general,
per-interface $\nunit^{\!\top}\Kmat\nunit$ is tightness for anisotropic
$\Kmat$ when the interface does \emph{not} happen to align with $\Kmat$'s
dominant eigendirection: in that case $\lambda_1(\Kmat)$ is still a safe
overestimate, just not an exact one. Given Corollary~\ref{cor:frobenius}
already bounds how loose that overestimate can be
($\leq\sqrt3\,\lambda_1$, only reached when $\Kmat$ is isotropic: in
which case $\lambda_1$ \emph{is} exact for every $\nunit$ regardless), this
is a good trade: it fixes both open gaps (the missing $\mu$ factor and the
Frobenius looseness) with a local, per-particle change, not an
architectural one.
\end{remark}

\subsection{The General Telegrapher Equation}
\label{sec:telegrapher}

Proposition~\ref{prop:main} was derived from the first-order system's
flux Jacobian. There is a second, classical route to the same physics:
eliminate $\Fdiff$ between Eqs.~\eqref{eq:mass_conservation}--\eqref{eq:cattaneo}
to get a single second-order-in-time equation (the Cattaneo/Maxwell-Cattaneo
route to the ``telegrapher's equation'', normally presented only for
scalar/isotropic diffusivity). This section derives it in two stages: an
\emph{exact} equation in the driver $\driver$ (Proposition~\ref{prop:telegrapher},
no approximation beyond what Eq.~\eqref{eq:cattaneo} already assumes), and
a \emph{local} specialisation to a closed equation in $U_0$ alone for any
general tensor $\Deff$ (Corollary~\ref{cor:telegrapher_local}), used as an
independent check on Proposition~\ref{prop:main} via a completely
different method (a dispersion relation, rather than a Jacobian
eigen-decomposition).

\begin{proposition}[General telegrapher equation, exact form]
\label{prop:telegrapher}
The conservation and Cattaneo-relaxation equations
\eqref{eq:mass_conservation}--\eqref{eq:cattaneo} imply, exactly (no
frozen-coefficient or isotropy assumption: only the two governing
equations themselves, and $\taurelax$ treated as locally constant so it
can be pulled through $\div$/$\partial_t$, exactly as in
Eq.~\eqref{eq:cattaneo} itself),
\begin{equation}
  \label{eq:telegrapher_q}
  \boxed{\;\taurelax\,\pdv[2]{U_0}{t} + \pdv{U_0}{t} = \div\!\big(\Kmat\grad \driver\big)\;} ,
\end{equation}
a PDE in the \emph{driver} $\driver$, not yet in $U_0$ alone.
\end{proposition}

\begin{proof}
Take $\partial_t$ of Eq.~\eqref{eq:mass_conservation}:
$\partial_t^2 U_0 + \div(\partial_t\Fdiff)=0$. Substitute
$\partial_t\Fdiff$ from Eq.~\eqref{eq:cattaneo},
\begin{equation}
  \partial_t\Fdiff = -\frac{\Kmat}{\taurelax}\grad \driver - \frac{\Fdiff}{\taurelax} ,
\end{equation}
giving
\begin{equation}
  \partial_t^2U_0 - \div\!\left(\frac{\Kmat}{\taurelax}\grad \driver\right) - \frac{1}{\taurelax}\div\Fdiff = 0 .
\end{equation}
Eliminate the remaining $\div\Fdiff$ using
Eq.~\eqref{eq:mass_conservation} itself, $\div\Fdiff=-\partial_tU_0$:
\begin{equation}
  \partial_t^2U_0 + \frac{1}{\taurelax}\partial_tU_0 - \div\!\left(\frac{\Kmat}{\taurelax}\grad \driver\right) = 0 .
\end{equation}
Multiplying by $\taurelax$ gives Eq.~\eqref{eq:telegrapher_q}.
\end{proof}

\begin{corollary}[Local form in $U_0$ alone]
\label{cor:telegrapher_local}
Under the \emph{additional} frozen-coefficient approximation of
Section~\ref{sec:quasilinear} ($\mu$, hence $\Deff=\mu\Kmat$, locally
constant in a neighbourhood, the same approximation already used to
derive the flux Jacobian, Section~\ref{sec:quasilinear}), $\grad\driver
\approx\mu\grad U_0$, and Eq.~\eqref{eq:telegrapher_q} reduces to
\begin{equation}
  \label{eq:telegrapher}
  \taurelax\,\pdv[2]{U_0}{t} + \pdv{U_0}{t} \approx \div\!\big(\Deff\grad U_0\big) ,
\end{equation}
now a closed equation in $U_0$ alone, for any symmetric PSD $\Deff$ (no
isotropy assumed). This is the form used for the local dispersion analysis
below: a plane-wave ansatz is itself only meaningful locally, over a
neighbourhood small enough that treating $\taurelax,\Deff$ as constant is
already implicit, so no additional assumption is introduced beyond what
the analysis needs anyway.
\end{corollary}
\emph{Eq.~\eqref{eq:telegrapher_q} is the one to trust globally}
($\driver$ is what the code actually evolves gradients of,
Section~\ref{sec:governing-system}, and it is exact, not an
approximation); Eq.~\eqref{eq:telegrapher} is a local specialisation,
useful specifically because it isolates $U_0$ for the dispersion relation
below, and reduces to Eq.~\eqref{eq:telegrapher_q} exactly whenever
$\driver=U_0$ (\code{isotropic\_constant}, $\mu=1$).

\subsubsection{Two Limits}
Eq.~\eqref{eq:telegrapher_q} interpolates between two regimes already
discussed separately elsewhere in this document, now visibly unified in
one equation:
\begin{itemize}
\item $\taurelax\to0$: the first term drops and
  Eq.~\eqref{eq:telegrapher_q} collapses to
  $\partial_tU_0=\div(\Kmat\grad \driver)$: Fick's law substituted into
  mass conservation, i.e.\ exactly the parabolic equation of
  Part~\ref{part:parabolic}, Eq.~\eqref{eq:parabolic_diffusion_rho_Z}.
  This limit is consistent with
  \code{chemistry\_part\_integrate\_flux\_source\_term}'s exact solution
  (\code{hyperbolic/chemistry\_flux.h}), which computes
  $\Fdiff_{\rm out}=\Fdiff_{\rm in}\exp(-\dt/\taurelax)+\Fdiff_{\rm
  parabolic}\big(1-\exp(-\dt/\taurelax)\big)$ and so also collapses onto
  $\Fdiff_{\rm out}\to\Fdiff_{\rm parabolic}$ as $\taurelax\to0$ for any
  $\dt$; the discretised flux update and the continuum PDE agree on
  this limit.
\item $\taurelax\to\infty$ (or, more precisely, on timescales
  $t\ll\taurelax$): the equation is dominated by
  $\taurelax\partial_t^2U_0\approx\div(\Kmat\grad\driver)$, an
  undamped wave equation, genuinely hyperbolic (finite-speed)
  propagation, the point of the whole regularisation.
\end{itemize}

\subsubsection{Dispersion Relation: An Independent Check on $\chyp$}
Seek plane-wave solutions $U_0=\hat U_0\exp[i(\mathbf k\cdot\mathbf x-\omega t)]$
of the \emph{local} form, Eq.~\eqref{eq:telegrapher}
(Corollary~\ref{cor:telegrapher_local}; frozen $\taurelax,\Deff$, exactly
what a plane-wave ansatz already requires to be self-consistent).
Substituting,
\begin{equation}
  -\taurelax\omega^2 U_0 - i\omega U_0 = -(\mathbf k^{\!\top}\Deff\,\mathbf k)\,U_0
  \quad\Longrightarrow\quad
  \taurelax\omega^2 + i\omega - \mathbf k^{\!\top}\Deff\,\mathbf k = 0 ,
\end{equation}
a quadratic in $\omega$ with roots
$\omega = \dfrac{-i \pm \sqrt{-1+4\taurelax\,\mathbf k^{\!\top}\Deff\mathbf k}}{2\taurelax}$.
Two limits of this dispersion relation reproduce results derived by
entirely different means elsewhere in this document:
\begin{itemize}
\item \textbf{Short wavelength} ($4\taurelax\,\mathbf k^{\!\top}\Deff\mathbf k\gg1$,
  i.e.\ well inside the propagating/hyperbolic regime): writing
  $\mathbf k=k\nunit$,
  \begin{equation}
    \mathrm{Re}(\omega) \approx \pm\sqrt{\frac{\mathbf k^{\!\top}\Deff\mathbf k}{\taurelax}}
    = \pm k\sqrt{\frac{\nunit^{\!\top}\Deff\nunit}{\taurelax}} ,
  \end{equation}
  so the phase speed $\mathrm{Re}(\omega)/k \to \pm\sqrt{\nunit^{\!\top}\Deff\nunit/\taurelax}$,
  exactly $\pm\chyp(\nunit)$ of Proposition~\ref{prop:main}, recovered
  here from a plane-wave dispersion analysis of the reduced scalar PDE
  instead of from the first-order system's flux-Jacobian eigenvalues. Two
  unrelated derivation methods agreeing is a meaningful independent check,
  not a restatement; Eq.~\eqref{eq:telegrapher} was derived by
  eliminating $\Fdiff$ algebraically, with no reference to Jacobians or
  eigenvectors. The (present in both derivations) $-i/(2\taurelax)$
  imaginary part of $\omega$ is a uniform decay rate: every propagating
  mode damps on the relaxation timescale $\taurelax$, consistent with
  $\Fdiff$ relaxing toward its Fickian value.
\item \textbf{Long wavelength} ($4\taurelax\,\mathbf k^{\!\top}\Deff\mathbf k\ll1$):
  $\omega\approx-i\,\mathbf k^{\!\top}\Deff\mathbf k$, purely diffusive
  decay with no oscillation, exactly the dispersion relation of the
  parabolic equation $\partial_tU_0=\div(\Deff\grad U_0)$ itself. This is
  the same $\taurelax\to0$ recovery as above, seen instead as a
  long-wavelength (rather than small-$\taurelax$) statement: even at
  fixed, finite $\taurelax$, sufficiently large-scale features of the
  solution always diffuse rather than propagate.
\end{itemize}

\section{Relaxation Time and Timestep}

\subsection{Physically-Motivated Normalisation of $\taurelax_0$}
\label{sec:tau0-physical}

The propositions above fix the \emph{shape} dependence of $\chyp$ but say
nothing about how large $\taurelax_0$ itself should be. Numerically,
$\taurelax_0$ is a free parameter-file constant; physically, in the
\code{turbulent\_mode} closure it is supposed to represent \emph{how many
local shear/turnover times the diffusive flux takes to relax to its
Fickian equilibrium}. This section derives why that constant cannot be
chosen independently of $\kappaC$, checks the resulting formula against an
independent literature derivation, and revises the ``does hyperbolic beat
parabolic'' criterion accordingly.

\subsubsection{Why $\taurelax_0$ Must Depend on $\kappaC$}
\label{sec:tau0-motivation}
Whether the hyperbolic scheme is faster than the parabolic one is decided
by comparing $\taurelax$ against the intrinsic parabolic crossing time of
the same problem, $\taurelax_{\rm par}\sim\dx^2/D$: exactly the ratio the
code already forms internally for the HLL blending factor
(\code{chemistry\_riemann\_compute\_hyperbolic\_blending\_factor},
\code{chemistry\_riemann\_utils.h}). Since $D$ (through $\kappaP$) is itself
set by $\kappaC$, this comparison timescale is not independent of
$\kappaC$: doubling $\kappaC$ halves $\taurelax_{\rm par}$, so a
$\taurelax_0$ that does not track $\kappaC$ silently drifts to a different
point on the parabolic$\leftrightarrow$hyperbolic spectrum every time
$\kappaC$ is retuned (e.g.\ during the flux-limiter calibration of
Section~\ref{sec:hyperbolic-limiter}). A physically meaningful
$\taurelax_0$ must therefore be defined \emph{relative to} $\kappaC$, not
alongside it.

\subsubsection{Should $\taurelax$ Use the Regularised $\Splus$?}
\label{sec:tau-no-regularisation}
Eq.~\eqref{eq:tau_turbulent} uses $\norm{\Shear}_F$, the Frobenius norm of
the \emph{unregularised} shear tensor, even though $\Kmat$ itself is built
from $\Splus$ in \code{anisotropic\_gradient} mode
(Section~\ref{sec:diffusion-tensor-recap}). This is deliberate, not an
inconsistency.

\begin{corollary}[The $\Splus/\Shear$ norm ratio is bounded]
\label{cor:shear-regularisation-bound}
For any nonzero traceless symmetric $3\times3$ tensor $\Shear$ with
eigenvalues $\lambda_1\geq\lambda_2\geq\lambda_3$ (tracelessness forces
$\lambda_1>0>\lambda_3$ whenever $\Shear\neq0$),
\begin{equation}
  \frac{\norm{\Splus}_F}{\norm{\Shear}_F} \in
  \left[\frac{1}{\sqrt3},\ \sqrt{\frac23}\,\right] \approx [0.577,\ 0.816] ,
\end{equation}
the lower end attained at $\Shear=\mathrm{diag}(1,1,-2)$ (up to scale/permutation),
the upper end at $\Shear=\mathrm{diag}(2,-1,-1)$.
\end{corollary}
\begin{proof}
$\norm{\Splus}_F^2=\sum_{\lambda_i>0}\lambda_i^2$. If $\lambda_2\leq0$
(one positive eigenvalue), write $\lambda_2=-t\lambda_1$,
$\lambda_3=-(1-t)\lambda_1$ for $t\in[0,\tfrac12]$ (tracelessness and
ordering); the squared ratio is
$1/[1+t^2+(1-t)^2]$, decreasing on $[0,\tfrac12]$ from $1$ (the excluded
$\Shear=0$ boundary) to $\tfrac23$ at $t=\tfrac12$
($\lambda_2=\lambda_3$, i.e.\ $\mathrm{diag}(2,-1,-1)$); this is the upper end.
If $\lambda_2>0$ (two positive eigenvalues), write $\lambda_2=t\lambda_1$,
$t\in(0,1]$, so $\lambda_3=-(1+t)\lambda_1$; the squared ratio is
$(1+t^2)/[2(1+t+t^2)]$, decreasing on $(0,1]$ from $\tfrac12$ (the
$\lambda_2\to0$ boundary shared with the other branch) to $\tfrac13$ at
$t=1$ ($\lambda_1=\lambda_2$, i.e.\ $\mathrm{diag}(1,1,-2)$); this is the lower
end.
\end{proof}

\begin{remark}[$\Shear$, not $\Splus$, is the right choice for $\taurelax$]
\label{rem:no-shear-regularisation-in-tau}
Two independent reasons. Physically, $\taurelax$ is a turbulent mixing
timescale: it must track the total local strain-rate magnitude, since both
the stretching and compressive legs of the flow contribute to how fast an
eddy folds and mixes a passive scalar, exactly the full-norm scaling in
Romano et al.'s Eq.~\eqref{eq:romano_tdiff} below. $\Splus$ is a different
kind of object: a device for keeping $\Kmat$ positive semi-definite so
that diffusion smooths gradients rather than sharpening them
(Section~\ref{sec:diffusion-tensor-recap}), with no independent standing
as a timescale, and a scalar $\taurelax$ cannot carry $\Splus$'s
directional information in any case; only $\Kmat$'s eigenvectors do that,
through $\Deff$ itself.

Quantitatively, Corollary~\ref{cor:shear-regularisation-bound} shows the
two choices would differ only by a factor in
$[\sqrt{3/2},\sqrt3]\approx[1.22,1.73]$, bounded, and fully degenerate
with $\taurelax_0$, which is already a free, per-simulation calibrated
constant. Regularising would not fix anything; it would only relabel
$\taurelax_0$, while making $\taurelax$ larger (since
$\norm{\Splus}_F\leq\norm{\Shear}_F$ always) and pushing the scheme
further from the parabolic limit, the opposite of the accuracy
requirement motivating this section (small $\taurelax$ recovers Fick's
law). $\norm{\Shear}_F$ is both the physically correct choice and the more
conservative one.
\end{remark}

\subsubsection{Cross-Check: The Parabolic Mixing Timescale of Romano et al.~(2022)}
Romano, Nagamine \& Hirashita~\citeyear{Romano2022} use the same
Shen-et-al.~\citeyear{ShenWadsleyStinson2010}/Smagorinsky-Lilly~\citeyear{Smagorinsky1963}
closure implemented here (their diffusivity, eq.~10 of that paper, is
$D=C_d\norm{\Shear_{ab}}h^2\rho$: identical in structure to our
$\kappaP$ for \code{isotropic\_smagorinsky}, up to the SPH kernel-support
factor $\gammak^2$). They derive the physical timescale for turbulent
diffusion to erase a contrast $\Delta A$ in some transported quantity $A$
(their eq.~13):
\begin{equation}
  \label{eq:romano_tdiff}
  \taurelax_{{\rm diff},A} \sim \left(\frac{A}{\Delta A}\right)
  \left(\frac{C_d}{0.01}\right)^{-1}
  \left(\frac{\norm{\Shear_{ab}}}{100\ {\rm km\,s^{-1}\,kpc^{-1}}}\right)^{-1}\ {\rm Gyr} ,
\end{equation}
i.e.\ $\taurelax_{{\rm diff},A}\propto(\kappaC\norm{\Shear}_F)^{-1}$. This
is precisely the $\kappaC^{-1}\norm{\Shear}_F^{-1}$ scaling
Section~\ref{sec:tau0-motivation} argues is required, an independent,
literature-derived confirmation that the coupling is physical, not merely
a numerical convenience. It also has no explicit dependence on resolution
$h$, for the same reason: $D\propto h^2$ cancels
the $h^2$ implicit in a diffusive crossing time $\dx^2/D$.

\subsubsection{The Coupled Closed Form}
Equation~\eqref{eq:tau_turbulent} states the adopted form,
$\taurelax=\taurelax_0/(\kappaC\norm{\Shear}_F)$. Propagating this through
$\chyp=\sqrt{\kappaP/(\rho\taurelax)}$ with the isotropic Smagorinsky
$\kappaP=\kappaC\gammak^2h^2\rho\norm{\Shear}_F$
(Corollary~\ref{cor:isotropic}) gives
\begin{equation}
  \label{eq:code_turbulent_new}
  \chyp^{\rm code} = \frac{\kappaC}{\sqrt{\taurelax_0}}\,\gammak\,h\,\norm{\Shear}_F ,
\end{equation}
i.e.\ $\kappaC$ appears to the first power, not under the square root:
coupling $\taurelax$ to $\kappaC$ contributes one factor of $\kappaC$
through $\kappaP\propto\kappaC$ and a second through
$1/\taurelax\propto\kappaC$, both inside $\chyp^2=\kappaP/(\rho\taurelax)$.
$\taurelax_0$ itself is dimensionless (a pure count of shear times),
since $\kappaC$ is dimensionless in the Smagorinsky closure; the
reparametrisation changes what value of $\taurelax_0$ corresponds to a
given physical relaxation time, but not its units.

\subsubsection{On Adding a $\driver/\Delta\driver$ Term}
\label{sec:no-q-over-dq}
A natural follow-up question: since Eq.~\eqref{eq:romano_tdiff} carries a
factor $A/\Delta A$, should $\taurelax$ carry an analogous
$\driver/\Delta\driver$ factor? \textbf{No.} $\taurelax_{{\rm diff},A}$
above is the timescale for the \emph{state} $A$ itself to relax
appreciably; it is a solution-dependent, emergent quantity that combines
the flux dynamics with how far the current state is from equilibrium.
$\taurelax$ in the Cattaneo system is a different object: a material/
turbulence property governing how fast the \emph{flux} responds to an
instantaneous local gradient, independent of how large the resulting
state change happens to be, exactly analogous to a mean-free time in the
kinetic-theory derivation of Fourier's law, which does not depend on the
temperature contrast either. The $\driver/\Delta\driver$-like behaviour is
not missing from this formulation: it emerges automatically as a
\emph{consequence} of solving Eqs.~\eqref{eq:mass_conservation}--\eqref{eq:cattaneo}
together (in the $\taurelax\to0$ limit the system converges to the
parabolic equation, whose own relaxation time is exactly
Eq.~\eqref{eq:romano_tdiff}). Building $\driver/\Delta\driver$ directly
into $\taurelax$ would double-count this and would additionally break the
frozen-coefficient assumption of Section~\ref{sec:quasilinear} ($\taurelax$
would depend on $U_0$ and its gradient, invalidating the flux-Jacobian
derivation of Proposition~\ref{prop:main} as stated, a real structural
cost, not just a modelling objection). Where a $\driver/\Delta\driver$-like
quantity legitimately belongs is in \emph{calibrating} the constant
$\taurelax_0$ itself: Eq.~\eqref{eq:romano_tdiff} can be evaluated at a
fiducial, order-unity contrast to translate a target physical mixing
timescale (in Gyr) into a concrete $\taurelax_0$, a calibration aid for
the constant, not a runtime dependence for $\taurelax$.

\subsubsection{Speed Criterion}
\label{sec:speed-criterion}
Using Eq.~\eqref{eq:code_turbulent_new} and $v_{\max}\approx2\chyp$
(comparable neighbouring particles, negligible relative velocity along
$\nunit$, the same simplifying assumption used throughout this
comparison):
\begin{equation}
  \dt_{\rm hyp} \approx \frac{\CCFL\dx}{2\chyp}
  = \frac{\CCFL\sqrt{\taurelax_0}}{2\,\kappaC\norm{\Shear}_F} ,
  \qquad
  \dt_{\rm par} = \frac{\CCFL}{\kappaC\norm{\Shear}_F}
\end{equation}
The second expression is the code's \code{isotropic\_smagorinsky}
parabolic timestep (\code{chemistry\_timesteps.h}), using
$\lambda_1(\Kmat)=\kappaP$, exact for isotropic $\Kmat$. The first uses
the identification $\gammak h=\dx$, matching the kernel support radius
to the Gizmo particle size up to an $O(1)$ geometric factor. The
resolution $h$, diffusivity normalisation
$\kappaC$, and local shear $\norm{\Shear}_F$ all cancel from the ratio:
\begin{equation}
  \label{eq:revised_ratio}
  \boxed{\;\frac{\dt_{\rm hyp}}{\dt_{\rm par}} = \frac{\sqrt{\taurelax_0}}{2} \;}
  \qquad\Longrightarrow\qquad
  \text{hyperbolic faster iff } \taurelax_0 \gtrsim 4 .
\end{equation}
Whether the hyperbolic scheme is faster than parabolic depends only on
the single dimensionless constant $\taurelax_0$, against a fixed
threshold of order unity.

\begin{remark}[Particle-size definitions]
\label{rem:particle-size-definitions}
The CFL length scale $\dx$ above is the same particle size used
throughout the parabolic scheme's own CFL bound, with the same choice of
GIZMO or Lanson \& Vila~\citeyear{LansonVila2008} definition and the same
caveat that switching between them at a fixed $\CCFL$ is not by itself a
route to a larger timestep (Part~\ref{part:parabolic},
Section~\ref{sec:particle-size}). Nothing about the hyperbolic timestep
changes that analysis; $\dx$ enters $\dt_{\rm hyp}$ exactly as it enters
$\dt_{\rm par}$ above.
\end{remark}

\subsubsection{Speed Criterion for \texttt{anisotropic\_gradient} Mode}
\label{sec:speed-criterion-aniso}
The derivation above is specific to \code{isotropic\_smagorinsky}: its
$\kappaP$ carries an explicit $\norm{\Shear}_F$ factor
(Section~\ref{sec:diffusion-tensor-recap}) that cancels the same factor
in $\taurelax$, and its parabolic timestep uses $\expr=\dx$
(\code{chemistry\_timesteps.h}), independent of the local
gradient. Neither holds for \code{anisotropic\_gradient}: its $\kappaP$
has no $\norm{\Shear}_F$ factor at all (shear enters only through
$\Splus$'s eigenvalues), and its parabolic timestep uses the
gradient-suppressed
$\expr=\norm{\driver}\dx/(\norm{\grad\driver}\dx+\norm{\driver})$
(\code{chemistry\_timesteps.h}).

Redo the ratio for this mode,
using $\lambda_1(\Splus)=\lambda_1(\Shear)$ (Corollary~\ref{cor:shear-regularisation-bound}'s
proof: the top eigenvalue of a nonzero traceless tensor is always
positive, hence unchanged by the $\Splus$ projection) and the same
$\gammak h=\dx$ identification as above:
\begin{align}
  \dt_{\rm hyp} &\approx \frac{\CCFL\dx}{2\chyp}
  = \frac{\CCFL\dx\sqrt{\taurelax_0}}{2\,\kappaC\gammak h\,\sqrt{\lambda_1(\Shear)\norm{\Shear}_F}}
  = \frac{\CCFL\sqrt{\taurelax_0}}{2\,\kappaC\sqrt{\lambda_1(\Shear)\norm{\Shear}_F}} , \\
  \dt_{\rm par} &= \frac{\CCFL\,\expr^2}{\kappaC\gammak^2h^2\lambda_1(\Shear)}
  = \frac{\CCFL\,\expr^2}{\kappaC\dx^2\lambda_1(\Shear)} ,
\end{align}
where the first line uses
$\chyp=\sqrt{\lambda_1(\Deff)/\taurelax}=\kappaC\gammak h\sqrt{\lambda_1(\Shear)\norm{\Shear}_F/\taurelax_0}$
(from $\Deff=\kappaC\gammak^2h^2\Splus$,
Section~\ref{sec:diffusion-tensor-recap}, and
Eq.~\eqref{eq:tau_turbulent}). Dividing:
\begin{equation}
  \label{eq:aniso_ratio}
  \frac{\dt_{\rm hyp}}{\dt_{\rm par}}
  = \frac{\sqrt{\taurelax_0}}{2}\,\left(\frac{\dx}{\expr}\right)^{\!2}
  \sqrt{\frac{\lambda_1(\Shear)}{\norm{\Shear}_F}} .
\end{equation}
This reproduces Eq.~\eqref{eq:revised_ratio} exactly when
$\expr=\dx$ (no gradient suppression) and $\Shear$ is isotropic
($\lambda_1(\Shear)/\norm{\Shear}_F=1$).

Write the local gradient steepness, in units of the resolution length,
\begin{equation}
  g \equiv \frac{\norm{\grad\driver}\,\dx}{\norm{\driver}} ,
\end{equation}
so that $g=0$ for a perfectly-resolved, flat field and $g\to\infty$ for a
steep, poorly-resolved one. In these terms, $\dx/\expr=1+g$.

Write the shear-anisotropy ratio
\begin{equation}
  r_{\Shear} \equiv \frac{\norm{\Shear}_F}{\lambda_1(\Shear)}
\end{equation}
(subscripted to avoid colliding with the unrelated ratio $r$ of
Section~\ref{sec:limiter-naive-port} or the radial coordinate $r$ of
Section~\ref{sec:exact-solutions}). Corollary~\ref{cor:shear-regularisation-bound}'s
proof bounds $\lambda_1(\Shear)/\norm{\Shear}_F\in[1/\sqrt3,\sqrt{2/3}\,]$
for any nonzero traceless $\Shear$ (same eigenvalue-triple argument, this
time for $\lambda_1$ rather than $\Splus$'s norm), i.e.
\begin{equation}
  r_{\Shear} \in [\sqrt{3/2},\sqrt6\,] \approx [1.22,\,2.45] ,
\end{equation}
tight at the same $\mathrm{diag}(2,-1,-1)$ (lower end) and
$\mathrm{diag}(1,1,-2)$ (upper end) triples.

Setting Eq.~\eqref{eq:aniso_ratio} to $1$:
\begin{equation}
  \label{eq:aniso_threshold}
  \boxed{\;\taurelax_0 \gtrsim \frac{4r_{\Shear}}{(1+g)^4}\;} ,
  \qquad r_{\Shear}\in[\sqrt{3/2},\sqrt6\,] .
\end{equation}
Unlike Eq.~\eqref{eq:revised_ratio}, this threshold is \emph{not} a
single universal constant: it depends on the local shear anisotropy
(through $r_{\Shear}$) and, more strongly (fourth power), on the local gradient
steepness $g$. Two limits are informative:
\begin{itemize}
\item \textbf{Smooth/well-resolved field} ($g\to0$, i.e.\
  $\expr\to\dx$, the same limit that makes
  \code{anisotropic\_gradient} degenerate to
  \code{isotropic\_smagorinsky}'s timestep form): the threshold reduces to
  $\taurelax_0\gtrsim4r_{\Shear}\in[4.9,\,9.8]$, \emph{higher} than the isotropic
  $4$, because $\kappaP$ lacks the $\norm{\Shear}_F$ boost
  \code{isotropic\_smagorinsky} gets, making $\Deff$ smaller for the same
  $\kappaC,h$; since $\dt_{\rm par}\propto1/\Deff$ grows faster than
  $\dt_{\rm hyp}\propto1/\sqrt{\Deff}$ as $\Deff$ shrinks, a smaller
  $\Deff$ favours parabolic, so hyperbolic needs a larger $\taurelax_0$ to
  compensate.
\item \textbf{Steep gradient} ($g\gg1$, under-resolved abundance
  contrasts, exactly the SN-injection-site regime motivating this whole
  investigation, Section~\ref{sec:scope}): the threshold falls as
  $g^{-4}$, well below $4$, so hyperbolic wins at \emph{smaller}
  $\taurelax_0$ than the isotropic estimate would suggest. This follows
  directly from $\dt_{\rm par}\propto\expr^2\propto(1+g)^{-2}$ shrinking
  with steepening gradients while $\dt_{\rm hyp}$ (a pure wave-speed CFL,
  Section~\ref{sec:tau0-motivation}) does not depend on $g$ at all.
\end{itemize}
The qualitative conclusion (hyperbolic's relative timestep advantage
\emph{grows} precisely where gradients are steepest) does not depend on
the $\gammak h=\dx$ identification: that identification enters
Eq.~\eqref{eq:aniso_ratio} only through the overall
$\sqrt{\taurelax_0}/2$ prefactor shared with the isotropic result, not
through the $g$- or $r_{\Shear}$-dependent factors, which are ratios of
same-mode quantities and are identification-independent.

\subsection{The Implemented Wave Speed}
\label{sec:confrontation}

\code{chemistry\_get\_physical\_hyperbolic\_soundspeed}
(\code{hyperbolic/chemistry\_getters.h}) is a single, general
formula, valid for every \code{diffusion\_mode} and
\code{relaxation\_time\_mode}:
\begin{equation}
  \label{eq:code_constant}
  \boxed{\;\chyp^{\rm code} = \sqrt{\frac{\lambda_1(\Deff)}{\taurelax}}\;} ,
  \qquad \Deff=\mu\Kmat\ \text{(\code{chemistry\_get\_physical\_matrix\_D})} ,
\end{equation}
exactly Eq.~\eqref{eq:main_result} of Proposition~\ref{prop:main} with
$\nunit^{\!\top}\Deff\nunit$ replaced by its particle-local, direction-free
bound $\lambda_1(\Deff)$ (Remark~\ref{rem:lambda-max-not-n}). No branch on
\code{relaxation\_time\_mode} is needed inside the getter: $\taurelax$ is
read from \code{p->chemistry\_data.tau}, already computed correctly for
whichever \code{relaxation\_time\_mode} is active by
\code{chemistry\_compute\_physical\_tau()} in
\code{chemistry\_prepare\_force()}, \emph{before} this getter is ever
called (from \code{chemistry\_iact.h} or the Riemann solver, both
during/after the force loop). So by the time $\chyp$ is needed, the
mode-dependence is already resolved upstream. Routing everything through
$\Deff$ (Section~\ref{sec:effective-diffusivity}) rather than carrying
$\Kmat$ and $\mu$ separately through the getter means there is only one
code path, with a single, unambiguous convention for $\mu$.

Clearing the speed threshold of
Sections~\ref{sec:speed-criterion}--\ref{sec:speed-criterion-aniso} is
not the end of the $\taurelax_0$ story: the same constant also controls
how the Riemann solver blends its Hopkins-limited and unlimited
dissipation branches, a tension worked out in
Section~\ref{sec:practical-tau0} once that solver has been introduced.

\section{Numerical Scheme}

\subsection{The Riemann Solver}
\label{sec:hyperbolic-riemann-solver}

Each interacting pair exchanges an interface flux by solving a Riemann
problem for the full four-component state $\Uvec=(U_0,\Fdiff)$
(\code{chemistry\_riemann\_solver\_HLL},
\code{hyperbolic/chemistry\_riemann\_HLL.h}), reusing the same face
geometry, i.e. area $\Aij$, position $\xij$, normal $\nunit$, as
Part~\ref{part:parabolic}, Section~\ref{sec:fvpm-geometry}. Unlike the
parabolic solver (Part~\ref{part:parabolic}, Section~\ref{sec:hll-flux}),
which has no wave speed of its own and borrows the hydrodynamic one,
this solver's wave speeds are built from $\chyp$ itself, in the same
PVRS (primitive-variable) form as the hydrodynamic estimate it replaces:
with $u_L,u_R$ the particle velocities projected onto $\nunit$ and
$\chyp^L,\chyp^R$ the two particles' own $\chyp$
(Section~\ref{sec:confrontation}),
\begin{equation}
  \label{eq:hyp_wavespeeds}
  \lambda_+ = \max(u_L+\chyp^L,\ u_R+\chyp^R) , \qquad
  \lambda_- = \min(u_L-\chyp^L,\ u_R-\chyp^R) ,
\end{equation}
guarded, like the parabolic solver, against near-degenerate
$\lambda_+-\lambda_-$ and pure-vacuum states, both returning a zero
flux.

The two standard HLL building blocks are then formed exactly as in
Part~\ref{part:parabolic}, Section~\ref{sec:hll-flux}, but now
component-wise over all four entries of $\Uvec$: with
$\mathrm{Flux}_L,\mathrm{Flux}_R$ the two particles' own flux tensors of
Section~\ref{sec:quasilinear}, projected onto $\nunit$, and
$U_L,U_R$ the reconstructed states of Section~\ref{sec:predictor},
\begin{equation}
  \label{eq:hyp_F_transport}
  F_{\rm transport} = \frac{\lambda_+\,\mathrm{Flux}_L - \lambda_-\,\mathrm{Flux}_R}{\lambda_+-\lambda_-} ,
  \qquad
  F_{\rm diss} = \frac{\lambda_+\lambda_-\,(U_R-U_L)}{\lambda_+-\lambda_-} ,
\end{equation}
each a four-component vector. Where the parabolic solver blends these
two building blocks with a minmod/artificial-diffusivity construction to
obtain its flux directly (Eq.~\eqref{eq:flux_hll}), this solver first
forms their unmodified sum, the plain, textbook HLL flux for the
hyperbolic system,
\begin{equation}
  \label{eq:hyp_F_HLL}
  F_{\rm HLL} = F_{\rm transport} + F_{\rm diss} .
\end{equation}
It is this \emph{raw} $F_{\rm HLL}$ that
Sections~\ref{sec:hll-consequence}--\ref{sec:hyperbolic-limiter} go on to
refine: they only ever rescale $F_{\rm diss}$ before it re-enters
Eq.~\eqref{eq:hyp_F_HLL}, never $F_{\rm transport}$.

\subsubsection{The Berthon Blend with the Parabolic Solver}
\label{sec:berthon-blend}
The pure hyperbolic flux $F_{\rm HLL}$ is not by itself a stable
discretisation of the $\taurelax\to0$ limit: as $\taurelax\to0$,
Eq.~\eqref{eq:hyp_wavespeeds}'s wave speeds diverge
($\chyp=\sqrt{\lambda_1(\Deff)/\taurelax}\to\infty$), taking the
CFL-stable but increasingly diffusive HLL dissipation with them, rather
than converging onto the accurate, TVD-limited parabolic flux of
Eq.~\eqref{eq:flux_hll}. Following Berthon
et~al.~\citeyear{Berthon2007}, the code instead blends $F_{\rm HLL}$ with
Part~\ref{part:parabolic}'s own Hopkins-HLL flux
(\code{chemistry\_riemann\_solver\_hopkins2017\_HLL}, evaluated on the
mass-density component alone using each particle's own local Fickian
flux $\Fdiff_{\rm parabolic}=-\Kmat\grad\driver$
(Remark~\ref{rem:cell-centred-gradients-only}) as its two one-sided
fluxes, giving a scalar $F_{\rm HLL}^{\rm par}$:
\begin{equation}
  \label{eq:berthon_blend}
  \boxed{\;
  F_0 = \alpha\, (F_{\rm HLL})_0 + (1-\alpha)\, F_{\rm HLL}^{\rm par} ,
  \qquad
  F_k = \alpha\, (F_{\rm HLL})_k \quad (k=1,2,3) \; ,}
\end{equation}
with $\alpha\in[0,1]$ a per-pair blending factor
(\code{chemistry\_riemann\_compute\_hyperbolic\_blending\_factor}),
$\alpha=\rho_\star/(0.1+\rho_\star)$ in terms of the pair's relaxation-time
ratio $\rho_\star$, given in full and motivated in
Eq.~\eqref{eq:berthon_alpha} (Section~\ref{sec:practical-tau0}). Only the mass-density
component $F_0$ has a parabolic counterpart to blend with: the
flux-vector components $F_{1,2,3}$ have no analogue in the parabolic
scheme (Part~\ref{part:parabolic} carries no $\Fdiff$ state to blend
against), so they are simply damped by $\alpha$ rather than blended.

As $\alpha\to0$ ($\taurelax\to0$, the near-equilibrium regime), $F_0\to
F_{\rm HLL}^{\rm par}$ exactly and $F_{1,2,3}\to0$: the mass flux
collapses onto the parabolic scheme's own accurate flux, and the
flux-vector components vanish, consistent with the $\taurelax\to0$
collapse of Eq.~\eqref{eq:cattaneo} onto Fick's law
(Section~\ref{sec:telegrapher}'s ``Two Limits''). As $\alpha\to1$ (the
wave-dominated regime), $F\to F_{\rm HLL}$ in full: the pure hyperbolic
HLL flux drives all four components, unmixed with any parabolic
contribution.

The resulting $F$ is scaled by the physical face area $\Aij$ and
accumulated into each particle's flux buffer using the same pairwise
sign convention as Part~\ref{part:parabolic},
Section~\ref{sec:time-integration}, subtracted from the left particle and
added to the right, extended component-by-component across all four
entries of $\Uvec$ rather than the single scalar $\rhoZ$.

\subsection{The Half-Timestep Predictor}
\label{sec:predictor}

\code{chemistry\_gradients\_predict} (\code{hyperbolic/chemistry\_gradients.h})
builds the left/right states handed to the Riemann solver by
extrapolating both the conserved density $U_0=\rhoZ$ and the flux
$\Fdiff$ from a particle's own position to the interparticle interface,
over a half-timestep $\dth=\tfrac12\min(\dt_i,\dt_j)$. The two extrapolations
differ in kind. $U_0$'s evolution equation, $\partial_tU_0=-\div\Fdiff$,
carries no source term, so the code uses the ordinary first-order
MUSCL-Hancock form, additive in space and time:
\begin{equation}
  U_0(\mathbf x_{\rm face},t^n+\dth) \approx
  U_0(\mathbf x_{\rm cell},t^n) + \grad U_0\cdot\Delta\mathbf x
  - \dth\,(\div\Fdiff)_{\mathbf x_{\rm cell}} ,
\end{equation}
with $\Delta\mathbf x=\mathbf x_{\rm face}-\mathbf x_{\rm cell}$ and
$\div\Fdiff$ itself built from the cell-centred flux gradient. This is
the same predictor form used for hydro's own MUSCL-Hancock step, and
needs no special treatment: $U_0$'s equation has no stiff source.

The spatial term $\grad U_0\cdot\Delta\mathbf x$ and the temporal term
$-\dth\,(\div\Fdiff)_{\mathbf x_{\rm cell}}$ are summed into a single
$\Delta U_0$ \emph{before} the extrapolated state is handed to
\code{chemistry\_slope\_limit\_face} (\code{hyperbolic/chemistry\_gradients.h},
the same call that limits $\Fdiff$ below). The face limiter
therefore bounds the full second-order (space-\emph{and}-time)
reconstruction of $U_0$, not just its spatial piece: the temporal term
carries no monotonicity guarantee of its own, so leaving it outside the
limiter's reach would let it push $U_0$ beyond its neighbours' range at a
sharp gradient.

$\Fdiff$'s equation, Eq.~\eqref{eq:cattaneo}, does have a stiff source
($-\Fdiff/\taurelax$), so a first-order-in-time Taylor step would be
inconsistent whenever $\dth\gtrsim\taurelax$. Therefore, the source term
is integrated exactly rather than linearised elsewhere in the scheme
(\code{chemistry\_part\_integrate\_flux\_source\_term}). The predictor
applies that same exact integrator, but with the two operations ordered
\emph{space, then time}:
\begin{equation}
  \label{eq:predictor-order}
  \Fdiff(\mathbf x_{\rm face},t^n+\dth) =
  \big[\Fdiff(\mathbf x_{\rm cell},t^n)+\grad\Fdiff\cdot\Delta\mathbf x\big]\,
  e^{-\dth/\taurelax} + \Fdiff_{\rm parabolic}(\mathbf x_{\rm cell})\,
  \big(1-e^{-\dth/\taurelax}\big)
\end{equation}
(\code{chemistry\_gradients.h}): first spatially reconstruct
$\Fdiff$ to the interface using its own stored gradient, then relax the
whole extrapolated value toward $\Fdiff_{\rm parabolic}$ exactly. Since $\Fdiff_{\rm parabolic}$
is only ever available at the cell centre
(Remark~\ref{rem:cell-centred-gradients-only}), the alternative order
(relax the cell-centred value first, then add the spatial gradient
afterwards) is also conceivable. The two are not equivalent, and, as
Remark~\ref{rem:reverse-order-wrong} below shows, the alternative is the
wrong choice.

Proposition~\ref{prop:predictor} makes precise what ``exact integrator''
means for the predictor: freezing $\Fdiff_{\rm parabolic}$ at its only
available cell-centre value turns $\Fdiff$'s governing equation into a
plain, per-point relaxation ODE with no spatial coupling, and
Eq.~\eqref{eq:predictor-order} is exactly that ODE's solution, not merely
a discretisation of it.
\begin{proposition}[Eq.~\eqref{eq:predictor-order} solves the frozen-$\Fdiff_{\rm parabolic}$ relaxation ODE exactly]
\label{prop:predictor}
Freeze $\Fdiff_{\rm parabolic}$ at its only available value,
$\Fdiff_{\rm parabolic}(\mathbf x_{\rm cell})$, constant over the cell
(the same approximation already forced by
Remark~\ref{rem:cell-centred-gradients-only}). Then
$\partial_t\Fdiff=-(\Fdiff-\Fdiff_{\rm parabolic})/\taurelax$ has no
spatial-derivative coupling left in it: it is an independent ODE at
every point $\mathbf x$. Given the linear initial profile
$\Fdiff(\mathbf x,t^n)=\Fdiff(\mathbf x_{\rm cell},t^n)+\grad\Fdiff\cdot(\mathbf x-\mathbf x_{\rm cell})$
supplied by the gradient estimator, its exact solution at every
$\mathbf x$ and every $t>t^n$ is
\begin{equation}
  \Fdiff(\mathbf x,t) = \Fdiff_{\rm parabolic} +
  \big[\Fdiff(\mathbf x,t^n)-\Fdiff_{\rm parabolic}\big]\,e^{-(t-t^n)/\taurelax} ,
\end{equation}
which at $\mathbf x=\mathbf x_{\rm face}$, $t=t^n+\dth$ is exactly
Eq.~\eqref{eq:predictor-order}.
\end{proposition}
\begin{proof}
Fix a point $\mathbf x$ and write $Y(t)=\Fdiff(\mathbf x,t)-\Fdiff_{\rm parabolic}$,
with $\Fdiff_{\rm parabolic}$ held constant in both space and time (the
freezing assumption above). Since the spatial coupling has already been
discarded, the governing equation reduces at this fixed $\mathbf x$ to
the linear, separable ODE $\dot Y=-Y/\taurelax$. Separating variables,
$d(\ln Y)/dt=-1/\taurelax$, integrates to
$Y(t)=Y(t^n)\,e^{-(t-t^n)/\taurelax}$. This derivation used no property
of $\mathbf x$, so it holds independently at every point; substituting
$Y(t^n)=\Fdiff(\mathbf x,t^n)-\Fdiff_{\rm parabolic}$ back in, and
evaluating at $\mathbf x=\mathbf x_{\rm face}$, $t=t^n+\dth$, gives the
proposition's claim.
\end{proof}

\begin{remark}[Why the reverse order is wrong]
\label{rem:reverse-order-wrong}
Decompose Eq.~\eqref{eq:predictor-order} into a cell-mean piece and a
gradient piece:
\begin{equation}
  \Fdiff(\mathbf x_{\rm face},t^n+\dth) =
  \underbrace{\Fdiff(\mathbf x_{\rm cell},t^n)e^{-\dth/\taurelax}
  +\Fdiff_{\rm parabolic}(1-e^{-\dth/\taurelax})}_{\text{exact cell-centre relaxation}}
  \;+\;
  \underbrace{\big(\grad\Fdiff\cdot\Delta\mathbf x\big)\,e^{-\dth/\taurelax}}_{\text{gradient, decayed}} .
\end{equation}
The gradient contribution decays at the same rate as the mean, since
$\Fdiff_{\rm parabolic}$ carries no gradient of its own in this local
approximation: both pieces of Proposition~\ref{prop:predictor}'s exact
solution decay identically. The reverse order (relax the cell-centred
value first, then
add $\grad\Fdiff\cdot\Delta\mathbf x$ \emph{un}decayed) would instead give
$\Fdiff_{\rm parabolic}+[\Fdiff(\mathbf x_{\rm cell},t^n)-\Fdiff_{\rm parabolic}]e^{-\dth/\taurelax}
+\grad\Fdiff\cdot\Delta\mathbf x$: as $\taurelax\to0$ this tends to
$\Fdiff_{\rm parabolic}+\grad\Fdiff\cdot\Delta\mathbf x$, a spurious residual
gradient inherited from the pre-relaxation, hyperbolic-regime flux field,
surviving even in the limit where the flux should have fully relaxed to
the smooth parabolic profile. Eq.~\eqref{eq:predictor-order}'s order
avoids this: as $\taurelax\to0$ it tends to $\Fdiff_{\rm parabolic}$ exactly, with no
residual gradient, matching the Cattaneo-to-Fick collapse of
Eq.~\eqref{eq:cattaneo} itself.
\end{remark}

\begin{remark}[The opposite limit, $\taurelax\to\infty$]
As $\taurelax\to\infty$, $e^{-\dth/\taurelax}\to1-\dth/\taurelax+O(\taurelax^{-2})$
and Eq.~\eqref{eq:predictor-order} tends, at leading order, to
$\Fdiff(\mathbf x_{\rm cell},t^n)+\grad\Fdiff\cdot\Delta\mathbf x$, a
pure spatial extrapolation with no time evolution, matching
Eq.~\eqref{eq:cattaneo}'s own source term vanishing in this limit. Both
limits behave as they should.
\end{remark}

The face slope limiter (\code{chemistry\_slope\_limit\_face},
\code{hyperbolic/chemistry\_slope\_limiter\_face.h}) is applied to
$\grad\Fdiff\cdot\Delta\mathbf x$ before the exact time integration in
\code{chemistry\_gradients.h}, consistent
with the space-then-time order of Proposition~\ref{prop:predictor}: it
bounds the purely spatial reconstruction that is then fed as the initial
condition into the exact relaxation solve, not a separate step operating
on an already time-evolved quantity.

The implemented order is not merely a reasonable choice between two
options: Proposition~\ref{prop:predictor} shows it is the only one
consistent with the exact solution of the locally-available (cell-centred
$\Fdiff_{\rm parabolic}$) relaxation problem.

\subsection{Consequence for the HLL Numerical Dissipation}
\label{sec:hll-consequence}

An overestimated $\chyp$ does not only affect the timestep; it directly
inflates $F_{\rm diss}$ of Eq.~\eqref{eq:hyp_F_transport}, the numerical
dissipation of the raw hyperbolic HLL flux, most visibly when $\alpha\to1$
(the large-$\taurelax$ regime of the Berthon blend,
Section~\ref{sec:berthon-blend}, where $F_{\rm diss}$ passes through to
the final flux essentially unmixed with the parabolic branch).
Specialising Eq.~\eqref{eq:hyp_F_transport} to the near-symmetric case
$\lambda_+=-\lambda_-\equiv\Lambda$ (a fair approximation whenever the
relative particle velocity along $\nunit$ is small compared to $\chyp$,
the diffusion-dominated regime this document is about) makes the effect
explicit:
\begin{equation}
  F_{\rm diss} = \frac{\lambda_+\lambda_-\,(U_R-U_L)}{\lambda_+-\lambda_-}
  = \frac{-\Lambda^2(U_R-U_L)}{2\Lambda} = -\frac{\Lambda}{2}(U_R-U_L) ,
\end{equation}
i.e.\ \emph{linear} in the wave-speed estimate $\Lambda$. Any factor by
which $\chyp$ is overestimated passes through directly into extra
numerical diffusion in the large-$\taurelax$ ($\alpha\approx1$) regime,
on top of, and independent from, a second lever: whether $F_{\rm diss}$
itself is limited in that regime the way Part~\ref{part:parabolic}'s
parabolic Hopkins-HLL solver limits its own dissipation
(Section~\ref{sec:hll-flux}). Since
$\chyp=\sqrt{\lambda_1(\Deff)/\taurelax}$ is exact for isotropic $\Deff$
and a tight, safe bound otherwise (Section~\ref{sec:confrontation}), the
wave-speed lever is already pulled as far as it goes; the flux-memory
limiter of Section~\ref{sec:hyperbolic-limiter} addresses the second.

\subsection{A Causal Bound on the Dissipation}
\label{sec:causal-bound}

The dissipation term itself, independent of any error in $\chyp$, has a
structural problem: $F_{\rm diss}=-\tfrac12\Lambda(U_R-U_L)$ is applied in
full across the entire pair separation
$r_{ij}=|\mathbf x_i-\mathbf x_j|$ (the same notation as
Section~\ref{sec:practical-tau0}) in a single step, regardless of how far
apart the two particles actually are. Read as a diffusivity, this hands
mass across $r_{ij}$ at rate $O(\chyp r_{ij})$, but the system's only
physical signal speed is $\chyp$ (Proposition~\ref{prop:main}): a
disturbance at one particle cannot have reached the other unless
$\chyp\dt_{ij}\gtrsim r_{ij}$, with $\dt_{ij}$ the pair's flux timestep
(the same pairwise minimum, $\min(\dt_i,\dt_j)$, entering the
half-timestep predictor of Section~\ref{sec:predictor}). Whenever
$\chyp\dt_{ij}<r_{ij}$, the HLL dissipation is therefore transporting
mass beyond the causal front the Cattaneo system allows: an $O(\chyp
r_{ij})$ upwind diffusivity that exceeds the physical
$\Deff\sim\chyp^2\taurelax$ (Eq.~\eqref{eq:main_result}) once
$r_{ij}\gtrsim\chyp\taurelax$, i.e.\ exactly the wave-dominated regime of
Section~\ref{sec:main-theorem}.

The fix multiplies all four components of $F_{\rm diss}$ by
\begin{equation}
  \label{eq:causal-bound}
  \boxed{\;\zeta_{ij} = \min\!\left(1,\ \frac{\chyp\,\dt_{ij}}{r_{ij}}\right)\;} ,
\end{equation}
before it is recombined with the transport term. $\zeta_{ij}=1$ leaves the
dissipation untouched whenever the wave has had time to cross the pair
within one step; otherwise $\zeta_{ij}<1$ rescales it down to the distance
actually reachable in $\dt_{ij}$. This converts the $O(\chyp r_{ij})$
upwind estimate into an $O(\chyp^2\dt_{ij})$ one, the same order as the
Lax-Wendroff (second-order, minimally dissipative) flux for a linear wave
equation, so dissipative transport can no longer outrun $\chyp$
regardless of $r_{ij}$. Every quantity in Eq.~\eqref{eq:causal-bound} is
already available at the point $F_{\rm diss}$ is assembled, so the bound
costs one division and one \code{min} beyond the existing dissipation
term.

Equation~\eqref{eq:causal-bound} is the hyperbolic solver's only cap on
$F_{\rm diss}$. The parabolic scheme bounds its own dissipation
differently, by a two-argument minmod between $(1+\psi)F_{\rm transport}$
and $F_{\rm transport}+F_{\rm diss}$, with
$\psi=$\code{chem\_data->hll\_riemann\_solver\_psi} a free,
per-simulation-tuned parameter (the Hopkins-HLL construction referenced in
Section~\ref{sec:hll-consequence}). Through the Berthon blend that
combines the two solvers, that minmod still acts on the parabolic
contribution to a hyperbolic run's flux, but not on the hyperbolic one.
The two mechanisms are not complementary: they rescale the same term in
pursuit of the same goal, suppressing the raw HLL dissipation's tendency
to over-transport, so applying both to the hyperbolic branch is
redundant, and at a tight $\psi$ the minmod binds first and leaves
$\zeta_{ij}$ inert. The causal bound is the preferable formulation of the
two, which is why it is the one the hyperbolic branch keeps: $\zeta_{ij}$
has no tunable parameter to calibrate per problem, and it is derived
directly from $\chyp$, the one quantity the wave-speed analysis of this
document (Proposition~\ref{prop:main}, Section~\ref{sec:confrontation})
already fixes as exact or a tight, safe bound, rather than from a free
constant chosen to match a given test problem.

\subsubsection{Practical Considerations}

\begin{itemize}
\item \textbf{Length scale.} The chemistry CFL condition ties $\dt$ to the
  local particle size over $\chyp$ (Section~\ref{sec:speed-criterion}), so
  $\chyp\dt_{ij}/r_{ij}$ is $O(\CCFL)$ for essentially every
  well-resolved pair, not a ratio that only binds in rare, badly-behaved
  configurations: Eq.~\eqref{eq:causal-bound} is a global reduction of the
  dissipation's reach, applied on (almost) every face every step. Because
  the von Neumann/Lax-Wendroff stability floor for this term is set by the
  particle size rather than by the pair separation $r_{ij}$ itself (a
  distinct, though comparable, length scale;
  Section~\ref{sec:practical-tau0}), an implementation may prefer the
  particle size in Eq.~\eqref{eq:causal-bound} in place of $r_{ij}$.
\item \textbf{A floor is required.} $\zeta_{ij}\to0$ exactly would zero
  out the dissipation entirely whenever $\chyp\to0$, which happens in
  \code{turbulent\_mode} wherever the local shear vanishes
  (Section~\ref{sec:tau-modes}), a legitimate part of the flow, not
  a pathology, and one where some numerical dissipation is still needed
  for stability. $\zeta_{ij}$ must therefore be floored at a small,
  strictly positive value rather than left to reach zero.
\item \textbf{Interaction with \code{hyperbolic\_limiter\_scope}.} Like
  the flux-memory limiter of Section~\ref{sec:hyperbolic-limiter}, the
  causal bound scales all four components of $F_{\rm diss}$, but whether
  components 1--3 (the flux-vector dissipation, as opposed to component 0,
  the mass-density dissipation) are actually rescaled should follow the
  same \code{hyperbolic\_limiter\_scope} choice the flux-memory limiter
  already respects: suppressing the flux-component dissipation
  unconditionally risks preserving, rather than dissipating, an
  already-stale flux memory.
\end{itemize}

\subsection{Practical Guidance for Choosing $\taurelax_0$}
\label{sec:practical-tau0}

Sections~\ref{sec:speed-criterion}--\ref{sec:speed-criterion-aniso} and
this one are not independent: the same $\taurelax_0$ that must clear the
speed threshold also fixes the blending factor $\alpha$
(\code{chemistry\_riemann\_compute\_hyperbolic\_blending\_factor},
\code{chemistry\_riemann\_utils.h}), which decides whether a
step runs through the Hopkins-limited parabolic branch ($\alpha\to0$) or
the plain, unlimited HLL branch of Section~\ref{sec:hll-consequence}
($\alpha\to1$).

\begin{equation}
  \label{eq:berthon_alpha}
  \alpha = \frac{\rho_\star}{0.1+\rho_\star} ,
  \qquad \rho_\star = \frac{\taurelax_\star}{\taurelax_{\rm num}} ,
\end{equation}
with $\taurelax_\star$ the harmonic mean of the pair's two relaxation times
(equal to $\taurelax/2$ when $\taurelax_L=\taurelax_R=\taurelax$), and
$\taurelax_{\rm num}=r_{ij}^2/D_\star$ the local parabolic crossing time
built from the actual pair separation $r_{ij}=|\mathbf x_i-\mathbf x_j|$
(distinct from the Gizmo particle size $\dx$ used in the speed-criterion
sections, a different, though comparable, length scale). For
\code{isotropic\_smagorinsky}+\code{turbulent\_mode}, substituting
$\taurelax=\taurelax_0/(\kappaC\norm{\Shear}_F)$ and
$\Deff=\kappaC\gammak^2h^2\norm{\Shear}_F$:
\begin{equation}
  \rho_\star = \taurelax_0\,\frac{\sqrt3\,\gammak^2h^2}{2\,r_{ij}^2} ,
\end{equation}
i.e.\ $\rho_\star$ is \emph{exactly} proportional to $\taurelax_0$, with a
fixed, resolution/shear/$\kappaC$-independent geometric prefactor. Since $\gammak h$
and $r_{ij}$ are comparable length scales ($\gammak>1$, and the kernel
support radius is typically a modest multiple of the interparticle
spacing, not a fraction of it), this prefactor is unlikely to be
$\ll1$, and $\alpha$ saturates fast regardless: $\alpha=0.5$ at
$\rho_\star=0.1$, $\alpha\approx0.91$ at $\rho_\star=1$,
$\alpha\approx0.98$ at $\rho_\star=4$. Since the speed-parity threshold
itself is $\taurelax_0\gtrsim4$ (isotropic) to $\gtrsim4$--$10$
(anisotropic, Eq.~\eqref{eq:aniso_threshold}'s $g=0$ limit), any
$\taurelax_0$ large enough to win on raw timestep speed is, in the
current implementation, already deep in the unlimited-HLL regime.

``Fast'' and ``Hopkins-limited/accurate'' are therefore not two
independently reachable regimes under $\taurelax_0$ tuning alone; they
are the same knob, pointing in opposite directions. The flux-memory
limiter of Section~\ref{sec:hyperbolic-limiter}, and the causal bound of
Section~\ref{sec:causal-bound}, each soften this tension from a
different angle.

\subsection{A Hopkins-2017-Style Limiter for the Hyperbolic Branch}
\label{sec:hyperbolic-limiter}

\subsubsection{Why the Obvious Port Doesn't Work}
\label{sec:limiter-naive-port}
The parabolic Hopkins-HLL solver limits its dissipation with
\code{chemistry\_riemann\_compute\_alpha} (\code{chemistry\_riemann\_utils.h}),
built around the ratio
\begin{equation}
  r = \frac{v_{\rm HLL}\,\dx\,\lvert\driver_\star\rvert}{\norm{\Kmat_\star}_F\,\lvert U_\star\rvert} ,
  \qquad v_{\rm HLL} = \tfrac12\lvert u_R-u_L\rvert + \tfrac12(c_{s,L}+c_{s,R}) ,
\end{equation}
comparing a numerical transport rate ($v_{\rm HLL}$, the \emph{hydro}
sound speed, an independent scale, since parabolic diffusion has no
wave speed of its own) against a physical diffusion rate (built from
$\Kmat$). The obvious port (reuse this formula verbatim, substituting
the genuine hyperbolic wave speed $\chyp$ for $v_{\rm HLL}$) does not
work: it is redundant with the blending factor $\alpha$
(Section~\ref{sec:practical-tau0}) already computed for the same pair.

\begin{remark}[The swapped ratio is a function of $\taurelax$ alone]
\label{rem:limiter-redundant}
Substituting $\chyp=\sqrt{\lambda_1(\Deff)/\taurelax}$ for $v_{\rm HLL}$
and $\driver_\star/U_\star=\mu$ (Corollary~\ref{cor:isotropic}) puts
$\Deff$ into both the numerator (through $\chyp$) and the denominator
(through $\Kmat_\star$) of $r$. Holding every other quantity fixed and
varying only $\taurelax$ (the one free parameter this section is about)
gives $r\propto\taurelax^{-1/2}$, while the existing blending ratio
$\rho_\star=\taurelax_\star/\taurelax_{\rm num}$
(Section~\ref{sec:practical-tau0}) is $\propto\taurelax^{+1}$: an exact
monotonic power-law relationship, $r\propto\rho_\star^{-1/2}$. The swapped $r$ therefore
carries no information beyond what $\alpha$ already encodes: both are
functions of the same scheme parameters ($\Deff$, $\taurelax$, $\dx$)
alone, since $\driver_\star/U_\star=\mu$ is a fixed per-mode constant,
not a diagnostic of the local metal field. A working limiter needs a
signal $\alpha$ does not already see: something that depends on the
actual local state, not just on scheme parameters.
\end{remark}

\subsubsection{Flux Memory}
\label{sec:limiter-flux-memory}
In the parabolic scheme, comparing dissipation against $\Kmat\grad\driver$
is exactly right because the flux \emph{is} $\Kmat\grad\driver$
algebraically; there is nothing else it could be. In the hyperbolic
scheme this is no longer true: $\Fdiff$ is a dynamical variable that only
relaxes toward $\Fdiff_{\rm parabolic}=-\Kmat\grad\driver$ on timescale
$\taurelax$ (Eq.~\eqref{eq:cattaneo}). An un-relaxed flux that persists
and keeps transporting metal after the local gradient that produced it
has moved on (not a wave-speed error) is the mechanism identified
(Section~\ref{sec:practical-tau0}) as the actual source of large-$\taurelax$
over-transport, and it is precisely what Remark~\ref{rem:limiter-redundant}'s
signal cannot see.

Define the \emph{flux memory} as the (pair-averaged, exactly as
$\Kmat_\star,\driver_\star,U_\star$ already are) residual between the
dynamical flux and its own equilibrium target,
\begin{equation}
  \Fdiff_{\rm memory} = \tfrac12(\Fdiff_L+\Fdiff_R)
  - \tfrac12(\Fdiff_{{\rm parabolic},L}+\Fdiff_{{\rm parabolic},R}) ,
\end{equation}
and the dimensionless ratio (subscripted to avoid colliding with the
unrelated $s\equiv\sqrt{t^2-r^2/\chyp^2}$ of Section~\ref{sec:exact-solutions})
\begin{equation}
  s_{\rm mem} = \frac{\norm{\Fdiff_{\rm memory}}}{\norm{\Fdiff_\star}+\norm{\Fdiff_{{\rm parabolic},\star}}} ,
  \qquad
  \boxed{\;\eta = \frac{1}{1+s_{\rm mem}}\;} ,
\end{equation}
multiplying $F_{\rm diss}$ (\code{chemistry\_riemann\_compute\_hyperbolic\_diffusivity\_limiter},
\code{chemistry\_riemann\_utils.h}). $\eta\to1$ (no suppression) when the
flux has caught up to its Fickian target; $\eta\to0$ (suppress) when it
is still coasting on a stale gradient. The sum-normalisation (rather than
normalising by $\Fdiff_{\rm parabolic}$ alone) avoids a $0/0$ blind spot
at a sharp initial discontinuity: there, each particle's own local
gradient (and hence $\Fdiff_{\rm parabolic}$) can be exactly zero
even though the pair straddles a real jump in $U_0$, which is legitimate
transport to kick-start, not numerical excess; the limiter correctly
leaves this case alone ($\eta=1$ when both $\Fdiff_\star$ and
$\Fdiff_{{\rm parabolic},\star}$ vanish). This functional form is a
starting point, not a derived result: Hopkins' own $r_{\rm term}$ is
similarly an empirically-tuned saturating function, not a first-principles
one.

\section{Validation and Verification}

\subsection{Exact Solutions of the Cattaneo Equation}
\label{sec:exact-solutions}

The scheme is validated against closed-form solutions of the equation it
discretises. Both the 1D and the 3D free-space Green's functions are
elementary; this section states them, because the structure of the 3D one
is not the naive analogue of the 1D one and the difference matters for how
a simulated profile should be read.

Throughout, the initial condition is the one the test problems use (a
metal-density spike with \emph{zero} initial flux,
$U_0(\mathbf x,0)=q_0\delta^d(\mathbf x)$, $\Fdiff(\mathbf x,0)=0$), and
$a\equiv1/(2\taurelax)$, $s\equiv\sqrt{t^2-r^2/\chyp^2}$.

\subsubsection{One Dimension}
\label{sec:exact-1d}

In 1D the solution has two pieces: a ballistic front carrying a finite
fraction of the mass at $|x|=\chyp t$, and a smooth interior behind it
\cite{MasoliverPorraWeiss1993} (restated in
\cite{PruestelMeierSchellersheim2013}, Eqs.~(2.1)--(2.5)),
\begin{equation}
  \label{eq:exact-green}
  U_0(x,t) = q_0e^{-at}\left[
    \underbrace{\tfrac12\delta\big(\chyp t-|x|\big)}_{\text{ballistic front}}
    +\underbrace{\Theta\big(\chyp t-|x|\big)\frac{a}{2\chyp}
      \left(I_0(as)+\frac{t\,I_1(as)}{s}\right)}_{\text{smooth interior}}
  \right] .
\end{equation}
Integrating each piece, the front carries a mass fraction $e^{-at}$ (split
evenly between $\pm\chyp t$) and the interior the remaining $1-e^{-at}$;
the two sum to $q_0$ identically.

The front is not a small correction. At $\taurelax=10$ it still holds
$97.5\%$ of the metal at $t=0.5$ and $77.9\%$ at $t=5$: for most of the
parameter range of interest the exact solution is dominated by a pair of
travelling spikes, not by the smooth interior.

\subsubsection{Three Dimensions}
\label{sec:exact-3d}

The 3D case is \emph{not} obtained by reusing Eq.~\eqref{eq:exact-green}.
It is still elementary, but both the front and the interior differ.

Substituting $U_0=e^{-at}w$ removes the first-order time derivative from
Eq.~\eqref{eq:telegrapher} and leaves a Klein-Gordon equation with a
positive mass-squared term,
\begin{equation}
  \label{eq:kg}
  \partial_t^2w = \chyp^2\nabla^2w + a^2w ,
\end{equation}
whose sign gives \emph{modified} Bessel functions rather than ordinary
ones. For radially symmetric fields $\nabla^2w=\tfrac1r\partial_r^2(rw)$,
so the 3D radial Green's function follows from the 1D one by
$w_3=-\tfrac{1}{2\pi r}\partial_rw_1$ with
$w_1=\tfrac{1}{2\chyp}\Theta(\chyp t-r)I_0(as)$. That propagator answers
the momentum-kick problem ($w=0$, $\partial_tw=\delta^3$); our initial
condition is the density spike ($w=\delta^3$, $\partial_tw=a\delta^3$,
the second following from $\partial_tU_0=0$ at $t=0$), so the physical
solution is $(\partial_t+a)$ applied to it. For the smooth interior $r<\chyp t$ this
gives
\begin{equation}
  \label{eq:U3d-smooth}
  U_0(r,t) = q_0e^{-at}\frac{a}{4\pi\chyp^3}\left[
    \frac{t\big(as\,I_0(as)-2I_1(as)\big)}{s^3}
    + \frac{a\,I_1(as)}{s}\right] ,
\end{equation}
which is strictly positive.

\begin{remark}[The 3D front carries three distributions, not one]
\label{rem:3d-front-dipole}
Differentiating $w_3$ produces a $\delta'(\chyp t-r)$ from the delta
itself, plus two further $\delta(\chyp t-r)$ terms: one from the $1/r$
geometry, one from $\partial_t\Theta(\chyp t-r)$. Collected as a radial
mass density (per unit $r$), the front is
\begin{equation}
  \label{eq:front-measure-3d}
  \mu_{\rm front}(r,t) = q_0e^{-at}\left[
    -\chyp t\,\delta'\!\big(r-\chyp t\big)
    +\Big(1+at+\tfrac12(at)^2\Big)\,\delta\big(r-\chyp t\big)\right] ,
\end{equation}
a signed dipole layer plus a monopole shell. The dipole carries no net
mass, so the 3D front mass fraction is
\begin{equation}
  \label{eq:front-mass-3d}
  \frac{M_{\rm front}(t)}{q_0} = e^{-at}\left(1+at+\tfrac12(at)^2\right) ,
\end{equation}
against $e^{-at}$ in 1D. Dropping either extra term breaks the mass
budget; with all three, interior plus front sum to $q_0$ to
$4\times10^{-16}$ over $\taurelax\in[10^{-1},10^2]$, $t\in[0.1,5]$. At
$\taurelax=10$ the 3D front holds $99.78\%$ of the metal at $t=5$ where
the 1D front holds $77.9\%$: in 3D the exact solution is, to excellent
approximation, a thin expanding shell around a nearly empty interior.

The dipole is genuine physics, not an artefact of the distributional
bookkeeping: it is the $\taurelax\to\infty$ (pure-wave) limit's Huygens
dipole layer, and it means the exact solution has a \emph{negative}
density lobe just inside the front. The Cattaneo equation does not
preserve positivity in three dimensions for sharp initial data; only the
smooth interior, Eq.~\eqref{eq:U3d-smooth}, is sign-definite.
\end{remark}

Because of the $\delta'$, Eqs.~\eqref{eq:U3d-smooth} and
\eqref{eq:front-mass-3d} should be evaluated directly rather than by
numerically inverting the Laplace-domain solution: Bromwich-type
quadrature rings at the front, and the closed form is both exact and
cheaper.

\subsubsection{Comparing Against a Simulated Profile}
\label{sec:exact-comparison}

Two practical consequences follow, both easy to get wrong.

\emph{A 3D run profiled along one axis must be compared with the
$y,z$-marginal of the 3D solution, which is exactly the 1D Green's
function.} Integrating Eq.~\eqref{eq:telegrapher} over the transverse
directions annihilates the transverse Laplacian, so the line density
$\lambda(x,t)=\iint U_0\,\dd y\,\dd z$ satisfies the 1D telegrapher
equation with a $\delta(x)$ initial condition: $\lambda$ \emph{is}
Eq.~\eqref{eq:exact-green}. Piece by piece,
\begin{equation}
  \label{eq:marginal-3d}
  \lambda(x,t) = 2\pi\!\int_{|x|}^{\chyp t}\! U_0(r,t)\,r\,\dd r
  \;+\; e^{-at}\,\frac{at+\tfrac12(at)^2}{2\chyp t}\,\Theta(\chyp t-|x|)
  \;+\; \frac{e^{-at}}{2}\Big[\delta(x-\chyp t)+\delta(x+\chyp t)\Big] ,
\end{equation}
and the first two terms reproduce the smooth interior of
Eq.~\eqref{eq:exact-green} to machine precision. Only the monopole and
$\partial_t\Theta$ parts of the front shell project flat; the $\delta'$
dipole projects onto the edge spikes at $x=\pm\chyp t$, which carry the
$e^{-at}$ front fraction familiar from 1D. A marginal drawn with the
whole front mass spread flat misplaces exactly that fraction, the
majority of all metal throughout the wave-dominated regime.
Note the identity does not run backwards: the 3D \emph{radial} profile
(Sec.~\ref{sec:exact-3d}) remains a different function, and the radial
and marginal views weight the front very differently.

\emph{Both solutions vanish identically outside $|x|>\chyp t$.} The front
sits exactly at the boundary and the interior carries an explicit
$\Theta$, so the exact solution has zero mass beyond the causal front in
any dimension. Any metal a simulation places there is numerical, which
makes $\int_{r>\chyp t+\epsilon_0}\rhoZ\,\dd V$ a reference-free
diagnostic: it depends only on the front's position, not on the shape or
amplitude of either curve.

\emph{Second moments must be read against the moment identity, and it
has a sting in 3D.} Multiplying Eq.~\eqref{eq:telegrapher} by $r^2$ and
integrating gives $\taurelax\ddot M_2+\dot M_2=2d\,\kappa$ in $d$
dimensions, hence exactly
\begin{equation}
  \label{eq:moment-identity}
  \langle r^2\rangle(t)
  = 2d\,\kappa\left[t-\taurelax\big(1-e^{-t/\taurelax}\big)\right] .
\end{equation}
In the wave regime this evaluates to $3\chyp^2t^2$ in 3D while the
support is only $r\le\chyp t$: the excess over the geometric maximum
$\chyp^2t^2$ is carried by the signed dipole in
Eq.~\eqref{eq:front-measure-3d}, whose second moment is
$2(\chyp t)^2e^{-at}$. A positivity-preserving scheme therefore
\emph{cannot} satisfy causality and Eq.~\eqref{eq:moment-identity}
simultaneously in the wave-dominated regime; the best a positive causal
solution can reach is $\langle r^2\rangle=(\chyp t+\epsilon_0)^2$, and
RMS comparisons should be read against that ceiling as well as against
Eq.~\eqref{eq:moment-identity}. Because the front carries most of the mass
in the wave-dominated regime (Section~\ref{sec:exact-3d}), $\langle
r^2\rangle$ must be accumulated as a genuine mass-weighted moment,
$\sum_i m_{Z,i}\,r_i^2/\sum_i m_{Z,i}$ over the particles' own positions
and metal masses; binning the metal density onto a fixed-width radial (or
Cartesian) grid first and computing the moment from the quantized profile
systematically under-counts the thin, fast-moving front shell whenever the
bin width is comparable to or larger than the front's radial extent,
biasing the diagnostic toward the smooth interior it is meant to be
checked against.

Finally, the seed in the test problems is a ball (or segment) of radius
$\epsilon_0$, not a point. Convolution adds variances for independent
zero-mean displacements, so a finite seed shifts the exact second moment
by $\tfrac35\epsilon_0^2$ in 3D and $\tfrac13\epsilon_0^2$ in 1D,
a correction that dominates at early times, when $\chyp t\lesssim\epsilon_0$.

\subsection{Conservation and Stability}
\label{sec:hyperbolic-validation}

\textbf{Conservation.} $U_0=\rhoZ$ obeys the same continuity form,
Eq.~\eqref{eq:mass_conservation}, as the parabolic scheme's own
Eq.~\eqref{eq:parabolic_diffusion_rho_Z}, exchanged across the same FVPM
interfaces by the same one-Riemann-solve-per-pair architecture
(Section~\ref{sec:scope}). Discrete conservation of total metal mass
therefore follows from the same pairwise-antisymmetric flux-update
argument as Part~\ref{part:parabolic}, Section~\ref{sec:time-integration}
(a shared flux, subtracted from one particle and added to the other,
cannot change their sum), and is not re-derived separately here. The
other three components of $\Uvec$, the flux $\Fdiff$ itself, are not
conserved quantities in this sense: they relax toward
$\Fdiff_{\rm parabolic}$ on timescale $\taurelax$
(Eq.~\eqref{eq:cattaneo}) and carry no analogous budget.

\textbf{Stability.} Proposition~\ref{prop:main} fixes the characteristic
speed $\chyp$ exactly (isotropic $\Deff$) or as a safe, tight bound
(Corollary~\ref{cor:frobenius}), so the wave-speed CFL condition of
Section~\ref{sec:speed-criterion} rests on solid ground. A
von-Neumann-style sufficient condition on $\CCFL$ for the full nonlinear
scheme remains open: the HLL flux, the causal bound of
Section~\ref{sec:causal-bound}, and the flux-memory limiter of
Section~\ref{sec:limiter-flux-memory} (itself flagged there as ``a
starting point, not a derived result'') act together in a way not closed
in closed form here, exactly as Part~\ref{part:parabolic},
Section~\ref{sec:timestep} leaves the fully nonlinear parabolic scheme's
stability unproven beyond its own linear, unlimited appendix bound
(Appendix~\ref{app:von-neumann}).

\subsection{Suggested Test Problems}
\begin{itemize}
  \item \textbf{Front-speed test}: seed an isolated metal peak and verify
        the numerical front tracks $\chyp t$ (Proposition~\ref{prop:main}),
        comparing against the 1D exact Green's function
        Eq.~\eqref{eq:exact-green} or its 3D counterpart
        Eq.~\eqref{eq:U3d-smooth}, on a uniform particle distribution.
  \item \textbf{Relaxation-limit recovery}: run the same peak at
        successively smaller $\taurelax_0$ (Eq.~\eqref{eq:revised_ratio})
        and verify the profile converges onto the parabolic
        \code{isotropic\_constant} solution of Part~\ref{part:parabolic},
        Eq.~\eqref{eq:pde_const}, the $\taurelax\to0$ limit of
        Eq.~\eqref{eq:telegrapher_q}.
  \item \textbf{3D front structure}: verify the simulated front carries
        the predicted signed monopole-plus-dipole mass fraction,
        Eq.~\eqref{eq:front-mass-3d}, rather than a purely positive
        shell, distinguishing genuine Cattaneo physics
        (Remark~\ref{rem:3d-front-dipole}) from numerical smoothing that
        would erase the dipole.
  \item \textbf{Causal-front diagnostic}: confirm no metal appears beyond
        $r>\chyp t+\epsilon_0$ (Section~\ref{sec:exact-comparison}), a
        reference-free check that depends only on the front's position,
        not on the shape or amplitude of either exact curve.
\end{itemize}


\appendix

% theory/ParabolicDiffusion/appendices.tex
% Shared appendices (mathematical derivations, nomenclature) for both
% the parabolic and hyperbolic parts. Input from
% GEAR_FVPM_Diffusion_Theory.tex (the master document); not standalone.

\section{Mathematical Derivations}

\subsection{Proof of Proposition~\ref{prop:gradient-consistency} (Linear Consistency of the Gradient Estimator)}
\label{app:gradient-consistency-proof}
For the assumed linear field, $Q_i - Q_j = \mathbf G\cdot(\mathbf x_i -
\mathbf x_j) = \mathbf G \cdot \mathbf{dx}_{ij}$. Substituting into
Eq.~\eqref{eq:gradient_estimator},
\begin{equation}
  (\grad Q)_i = \sum_j (\mathbf G\cdot\mathbf{dx}_{ij})\, W_{ij}\, \mathbf B_i \cdot \mathbf{dx}_{ij}
  = \mathbf B_i \left(\sum_j W_{ij}\, \mathbf{dx}_{ij}\otimes\mathbf{dx}_{ij}\right) \mathbf G ,
\end{equation}
using the identity $(\mathbf a\cdot\mathbf u)(\mathbf M\mathbf u) =
\mathbf M(\mathbf u\otimes\mathbf u)\mathbf a$ for a vector $\mathbf a$,
matrix $\mathbf M$ and vector $\mathbf u$, applied with $\mathbf a =
\mathbf G$, $\mathbf u=\mathbf{dx}_{ij}$ componentwise, and factoring
$\mathbf B_i$ (constant over the sum) out to the left. The sum in
parentheses is exactly $\mathbf E_i$ of Eq.~\eqref{eq:E_matrix_def}: the
gradient estimator (Eq.~\eqref{eq:gradient_estimator}) and the moment
matrix are built from the identical per-neighbour weight $W_{ij}$, so no
further correction factor is needed. Hence
$(\grad Q)_i = \mathbf B_i\, \mathbf E_i\, \mathbf G = \mathbf E_i^{-1}\mathbf E_i\, \mathbf G = \mathbf G$,
since $\mathbf B_i \equiv \mathbf E_i^{-1}$ by definition.

\subsection{From Conservation and Fick's Law to the Parabolic Equation}
\label{app:derivation-pde}
Starting from the conservation law for a generic scalar density
$\rho Q$ advected with the flow and subject to a diffusive current
$\Fdiff$,
\begin{equation}
  \pdv{(\rho Q)}{t} + \div(\rho Q\,\mathbf v) = -\div\Fdiff ,
\end{equation}
the Lagrangian (particle-following) derivative satisfies
$\dv{}{t}(\rho Q) \equiv \pdv{(\rho Q)}{t} + \div(\rho Q\,\mathbf v) - \rho Q\,\div\mathbf v$
only up to the usual continuity-equation bookkeeping; in the FVPM
formalism this is sidestepped entirely because the conserved variable is
carried \emph{by} the moving, mass-conserving particle rather than
evolved on a fixed grid, so the advective flux never needs to be
discretised for a passive scalar: only the diffusive part remains to be
integrated,
\begin{equation}
  \pdv{(\rho Q)}{t}\bigg|_{\rm particle} = -\div\Fdiff .
\end{equation}
Substituting the tensorial Fick's law $\Fdiff = -\Kmat\cdot\grad q$
(Eq.~\eqref{eq:ficks_law}) gives the general parabolic equation
\begin{equation}
  \pdv{(\rho Q)}{t} = \div\big(\Kmat\cdot\grad q\big) ,
\end{equation}
which specialises to Eqs.~\eqref{eq:pde_const}--\eqref{eq:pde_aniso}
according to whether $q=\rho Q$ or $q=Q$ and whether $\Kmat$ is a scalar
multiple of the identity or the full anisotropic tensor $\kappaP\Splus$.

\subsection{Von Neumann Stability of the Explicit Parabolic Update}
\label{app:von-neumann}
Consider the 1D, constant-coefficient model problem
$\partial_t u = \kappaC\,\partial_x^2 u$ discretised with forward Euler
in time and central differences in space on a uniform grid of spacing
$\dx$ (the FVPM scheme reduces to this discretisation locally, in a
region of uniform particle spacing and constant $\Kmat=\kappaC\mathbf I$,
which is precisely the regime the code's
Eq.~\eqref{eq:dt_const} targets):
\begin{equation}
  \frac{u_i^{n+1}-u_i^n}{\dt} = \kappaC\,\frac{u_{i+1}^n - 2u_i^n + u_{i-1}^n}{\dx^2} .
\end{equation}
\begin{proposition}[Parabolic CFL bound]
The scheme above is Von Neumann stable if and only if
$\dt \le \dfrac{\dx^2}{2\kappaC}$.
\end{proposition}
\begin{proof}
Insert a Fourier mode $u_i^n = \hat u^n e^{\mathrm i k x_i}$, $x_i = i\,\dx$:
\begin{equation}
  \hat u^{n+1} = \hat u^n\Big[1 + \frac{\kappaC\dt}{\dx^2}\big(e^{\mathrm i k\dx} - 2 + e^{-\mathrm i k\dx}\big)\Big]
  = \hat u^n\Big[1 - \frac{4\kappaC\dt}{\dx^2}\sin^2\!\Big(\frac{k\dx}{2}\Big)\Big] ,
\end{equation}
using $e^{i\theta}+e^{-i\theta}-2 = -4\sin^2(\theta/2)$. The amplification
factor is
$G(k) = 1 - \dfrac{4\kappaC\dt}{\dx^2}\sin^2\!\big(\tfrac{k\dx}{2}\big) \in
\Big[1-\dfrac{4\kappaC\dt}{\dx^2},\ 1\Big]$
as $\sin^2(\cdot)$ ranges over $[0,1]$. Stability requires
$|G(k)|\le 1$ for every admissible $k$; since $G(k)\le1$ always, the
binding constraint is the lower bound,
\begin{equation}
  1 - \frac{4\kappaC\dt}{\dx^2} \ge -1
  \quad\Longleftrightarrow\quad
  \dt \le \frac{\dx^2}{2\kappaC} ,
\end{equation}
attained in the worst case $k\dx=\pi$ (the grid (particle)-scale
Nyquist mode). This is both necessary (violating it makes $|G(\pi/\dx)|>1$,
so that mode grows without bound under repeated application) and
sufficient (it makes $|G(k)|\le1$ for every $k$, so no mode grows).
\end{proof}
The code's constant $\CCFL$ (\code{chem\_data->C\_CFL\_chemistry}) absorbs
the factor of $1/2$ above together with an additional safety margin for
the (nonlinear, multi-dimensional, variable-coefficient, explicitly
minmod-limited) production scheme, for which the linear analysis above
is necessary but not sufficient. In the anisotropic case, replacing
$\kappaC \to \norm{\Kmat}_F$ in the bound is the standard generalisation
used when the exact eigenvalue decomposition of $\Kmat$ is not tracked
per-timestep for performance reasons: since
$\norm{\Kmat}_F \ge \lambda_{\max}(\Kmat)$ for any symmetric PSD matrix
(the Frobenius norm bounds the spectral radius from above), using
$\norm{\Kmat}_F$ in place of the true largest eigenvalue in
Eq.~\eqref{eq:dt_const}/\eqref{eq:dt_aniso} yields a \emph{more}
restrictive (i.e.\ still stable, possibly conservative) timestep bound.

\section{Nomenclature}

\subsection{Symbol to Code-Identifier Map}
\begin{center}
\begin{tabular}{@{}lll@{}}
\toprule
Symbol & Code identifier & Meaning \\
\midrule
$\rho$ & \code{hydro\_get\_physical\_density(p)} & Physical gas density \\
$\Zfrac$ & \code{chemistry\_get\_metal\_mass\_fraction} & Metal mass fraction \\
$\rhoZ$ & \code{chemistry\_get\_physical\_metal\_density} & Physical metal mass density \\
$m_Z$ & \code{p->chemistry\_data.metal\_mass[m]} & Conserved metal mass \\
$\driver$ & n/a (selected by mode) & Diffusion driver: $\rhoZ$ or $\Zfrac$ \\
$\kappaC$ & \code{chem\_data->diffusion\_coefficient} & Global diffusion normalisation \\
$\kappaP$ & \code{p->chemistry\_data.kappa} & Per-particle diffusion coefficient \\
$\Kmat$ & \code{chemistry\_get\_physical\_matrix\_K} & Diffusion tensor \\
$\Kstar$ & \code{chemistry\_riemann\_compute\_K\_star} & Interface-averaged diffusion tensor \\
$\Shear$ & \code{chemistry\_get\_physical\_shear\_tensor} & Traceless velocity shear tensor \\
$\Splus$ & \code{chemistry\_regularize\_shear\_tensor} & PSD-regularised shear tensor \\
$\Fdiff$ & \code{chemistry\_get\_physical\_parabolic\_flux} & Diffusive flux \\
$\gammak$ & \code{kernel\_gamma} & Kernel compact-support factor \\
$\mathbf B_i$ & \code{p->geometry.matrix\_E} (already inverted) & Least-squares gradient matrix \\
$V_i$ & \code{p->geometry.volume} & Particle effective volume \\
$\mathbf A_{ij}$ & (local, \code{chemistry\_iact.h}) & Generalised interface area vector \\
$\xij$ & (local, \code{chemistry\_iact.h}) & Interface position relative to $i$ \\
$\psiHLL$ & \code{chem\_data->hll\_riemann\_solver\_psi} & HLL transport-term weight \\
$\epsHLL$ & \code{chem\_data->hll\_riemann\_solver\_epsilon} & Direct-flux sign-consistency tolerance \\
$\alpha$ & \code{chemistry\_riemann\_compute\_alpha} & Artificial diffusivity \\
$\CCFL$ & \code{chem\_data->C\_CFL\_chemistry} & Chemistry CFL coefficient \\
$\dx$ & (local, \code{chemistry\_timesteps.h}) & GIZMO physical particle size \\
\midrule
\multicolumn{3}{@{}l}{\emph{Hyperbolic (Part~\ref{part:hyperbolic}) only}} \\
\midrule
$\taurelax$ & \code{p->chemistry\_data.tau} & Relaxation time \\
$\chyp$ & \code{chemistry\_get\_physical\_hyperbolic\_soundspeed} & Hyperbolic characteristic (wave) speed \\
$\Deff$ & \code{chemistry\_get\_physical\_matrix\_D} & Effective diffusivity, $\Deff=\Kmat(\driver/U_0)$ \\
$\Uvec$ & (assembled, \code{hyperbolic/chemistry\_getters.h}) & Conserved state vector $(\rhoZ,F_x,F_y,F_z)$ \\
$\Gflux{k}$ & (assembled, \code{hyperbolic/chemistry\_getters.h}) & Flux tensor, direction-$k$ component \\
$\expr$ & (local, \code{chemistry\_timesteps.h}) & Gradient-suppression effective length (parabolic's $E$, Eq.~\eqref{eq:E-bracket}) \\
\bottomrule
\end{tabular}
\end{center}

\subsection{Subscripts and Superscripts}
\begin{itemize}[noitemsep,topsep=0pt]
  \item $i,j$: particle indices (also used as tensor component indices where unambiguous from context)
  \item $L,R$: left/right Riemann states
  \item $n$: explicit timestep index
  \item $\star$: interface-averaged quantity
  \item $+$: positive-eigenvalue-projected (regularised) tensor
\end{itemize}

\subsection{Abbreviations}
\begin{itemize}[noitemsep,topsep=0pt]
  \item FVPM: Finite Volume Particle Method
  \item MFM: Meshless Finite Mass (GIZMO family of schemes)
  \item HLL: Harten--Lax--van Leer (approximate Riemann solver)
  \item CFL: Courant--Friedrichs--Lewy
  \item LES: Large-Eddy Simulation
  \item PSD: Positive Semi-Definite
\end{itemize}


\bibliography{references}

\end{document}
